\documentclass[11pt]{article}
\usepackage[T1]{fontenc}
\usepackage[utf8]{inputenc}
\usepackage{lmodern,amsmath,amssymb,amsthm,mathtools,microtype}
\usepackage[margin=28mm]{geometry}
\usepackage[authoryear,round]{natbib}
\usepackage{xurl,xcolor}
\usepackage[colorlinks=true,allcolors=blue!50!black]{hyperref}
\newtheorem{theorem}{Theorem}[section]
\newtheorem{proposition}[theorem]{Proposition}
\newtheorem{lemma}[theorem]{Lemma}
\newcommand{\workbenchurl}{https://github.com/KokunoYumeto/schemes-and-moduli-workbench}
\newcommand{\recordurl}{https://github.com/KokunoYumeto/mathematics-workbenches}
\providecommand{\tightlist}{\setlength{\itemsep}{0pt}\setlength{\parskip}{0pt}}
\title{Sharp residues and hypersurface monodromy, after Grothendieck and Katz}
\author{}
\date{23 November 2026}
\begin{document}
\maketitle
\begin{abstract}
Grothendieck's crystalline notes ask whether hypersurfaces in every nonempty parameter open realize each point-count residue modulo the characteristic after finite extension of the ground field. For degree-$d$ hypersurfaces in $\mathbf P^n$, $n\ge2$, we prove the exact threshold $d=n+2$ and determine the finite monodromy image in every characteristic and degree. At and above the threshold the geometric and arithmetic images are the full general linear group, except when the characteristic is two and $d\equiv n+1\pmod2$; there they are the full stabilizer of a canonical nonzero covector. The lower degrees have zero or one-dimensional fixed-vector spaces with their exact images. A cubic-and-hyperplane boundary has full generic Hasse--Witt rank and independently supported elliptic and three-point monodromy. Its traces, combined with a finite-cover twist, yield every residue in every parameter open. Moving divisors prove the plane groups. In higher dimension, positioned homologies give the groups for $p\ge5$, a modular norm calculation forces transvections for $p=3$, and an explicit elliptic double cover supplies the required extension transformations for $p=2$. Elementary coordinate operations finish the small-characteristic cases. Generic invertibility is used in Adolphson--Sperber's combinatorial form, and Katz's coefficient and point-count formulas are credited precisely. The mathematics and text were produced and checked by GPT-6 Astra (ChatGPT Pro mode); no human has checked the mathematics.
\end{abstract}

\section{The sharp threshold and the monodromy statements}
Grothendieck asks for a smooth section with cardinality ``congru mod $p$ \`a un \'el\'ement de $\underline F_p$ donn\'e \`a l'avance'' in cote~4, \S8.6, archive sheet~11, PDF image~12 \citep{Grothendieck}. The question quantifies over every nonempty parameter open and permits a finite extension of the finite ground field. For projective hypersurfaces the answer has an exact degree threshold.

Fix $n\ge2$, a prime $p$, and $d\ge1$. Let $S_d$ be the projective space of degree-$d$ equations in $n+1$ variables. A hypersurface is \emph{Hasse--Witt ordinary} here if Frobenius on $H^{n-1}(Y,\mathcal O_Y)$ is bijective; no stronger Hodge--Witt ordinarity is assumed. On the smooth ordinary open, its fixed vectors form a finite \'etale $\mathbf F_p$ local system $V$ of rank
\[
h=\binom{d-1}{n},\qquad
I_d=\{a\in\mathbf Z_{>0}^{n+1}:|a|=d\},
\]
where $h=0$ for $d\le n$.

\begin{theorem}[Sharp point-count residue threshold]\label{thm:threshold}
The smooth Hasse--Witt ordinary open is nonempty for every $p,n,d$ above. If $d\ge n+2$, the geometric monodromy image on $V$ has every element of $\mathbf F_p$ as a trace. Consequently, for every finite field $k$ of characteristic $p$, every nonempty open $U\subset S_d$ defined over $k$, and every $a\in\mathbf F_p$, there are a finite extension $k'/k$ and a smooth hypersurface $Y\in U(k')$ such that
\[
\#Y(k')\equiv a\pmod p.
\]
The threshold is sharp. If $d\le n$, every degree-$d$ hypersurface over every finite extension has residue $1$. If $d=n+1$, residue $1$ is absent on the nonempty smooth ordinary open over every finite extension.
\end{theorem}

The geometric ingredient is a joint-image statement for low-dimensional sections. Let $\kappa$ be algebraically closed of characteristic $p$, and fix $m\ge n-1$ hyperplanes $H_1,\ldots,H_m$ in general position in $\mathbf P^n_\kappa$. For index sets of the indicated sizes put
\[
\Lambda_I=\bigcap_{i\in I}H_i\cong\mathbf P^2\quad(|I|=n-2),
\qquad
\ell_J=\bigcap_{j\in J}H_j\cong\mathbf P^1\quad(|J|=n-1).
\]
General position means that an intersection of $r\le n$ distinct hyperplanes has codimension $r$, and an intersection of $n+1$ is empty. For $n=2$, the empty set $I$ means the ambient plane. Let $S$ be the cubic parameter open on which $Q$ is smooth and transverse to every positive-dimensional hyperplane stratum, avoids all zero-dimensional strata, and every plane cubic $E_I=Q\cap\Lambda_I$ is smooth ordinary. Write $Z_J=Q\cap\ell_J$.

\begin{theorem}[Independent cubic-boundary monodromy]\label{thm:joint}
The open $S$ is nonempty. Each $Z_J\to S$ is finite \'etale of degree $3$, and each $H^1_{\mathrm{et}}(E_I,\mathbf F_p)$ has rank $1$. Their joint geometric monodromy image is
\[
\prod_{|J|=n-1}S_3\ \times\
\prod_{|I|=n-2}\mathbf F_p^\times.
\]
In particular each factor group occurs supported on its own factor and as the identity on every other listed factor. The elliptic factor groups are trivial when $p=2$.
\end{theorem}

The following results determine every finite monodromy image and the exact canonical-covector obstruction in characteristic two.

\begin{theorem}[Exact images and an invariant covector]\label{thm:images}
Over every perfect ground field of characteristic $p$, the geometric and arithmetic monodromy images on $V$ are trivial if $d\le n$, and equal to $\mathrm{GL}_1(\mathbf F_p)$ if $d=n+1$. If $p=2$ and $d=n+2$, both images are $\mathrm{GL}_{n+1}(\mathbf F_2)$. If $p=2$ and $d=n+1+2e$ with $e\ge0$, there is a canonical nowhere-zero morphism of local systems $\lambda:V\to\mathbf F_2$. For $e\ge1$ its kernel is a proper nonzero invariant hyperplane. Both images lie in its stabilizer. Equality with that stabilizer is determined in Theorems~\ref{thm:plane-images} and~\ref{thm:higher-images}.
\end{theorem}

\begin{proposition}[Determinants and rank-one semisimple elements]\label{prop:determinant-reflections}
Let $d\ge n+1$ and $h=\binom{d-1}{n}$. Over every perfect ground field of characteristic $p$, the geometric and arithmetic images on $V$ have determinant image $\mathbf F_p^\times$. If $p$ is odd and $d\ge n+2$, then for each $a\in\mathbf F_p^\times\setminus\{1\}$ the geometric image contains an element conjugate over $\mathbf F_p$ to
\[
\operatorname{diag}(a,1,\ldots,1).
\]
Thus it contains a semisimple element whose difference from the identity has rank $1$, with any prescribed nontrivial eigenvalue in $\mathbf F_p^\times$.
\end{proposition}

\begin{theorem}[Full plane-curve monodromy]\label{thm:plane-images}
Let $n=2$, $d\ge4$, and $g=(d-1)(d-2)/2$. On the smooth Hasse--Witt ordinary degree-$d$ plane-curve parameter open over every perfect ground field of characteristic $p$, the geometric and arithmetic monodromy images on $V$ coincide. If $p$ is odd, or if $p=2$ and $d$ is even, they are
\[
\mathrm{GL}_{g}(\mathbf F_p).
\]
If $p=2$ and $d\ge5$ is odd, let $\lambda:V\to\mathbf F_2$ be the canonical nonzero covector in Theorem~\ref{thm:images}. Both images are its full stabilizer:
\[
\operatorname{Stab}(\lambda)
\simeq \mathbf F_2^{\,g-1}\rtimes\mathrm{GL}_{g-1}(\mathbf F_2).
\]
\end{theorem}

\begin{theorem}[All higher-dimensional images]\label{thm:higher-images}
Let $n\ge3$, $d\ge n+2$, and let the ground field be perfect of characteristic $p$. Put $h=\binom{d-1}{n}$. The geometric and arithmetic images on the smooth Hasse--Witt ordinary degree-$d$ hypersurface parameter open coincide and are
\[
\begin{cases}
\mathrm{GL}_h(\mathbf F_p),&p\text{ odd},\\
\mathrm{GL}_h(\mathbf F_2),&p=2,\ d\equiv n\pmod2,\\
\operatorname{Stab}(\lambda),&p=2,\ d\equiv n+1\pmod2,
\end{cases}
\]
where $\lambda$ is the canonical nonzero covector of Theorem~\ref{thm:images}. In the last case
\[
\operatorname{Stab}(\lambda)\simeq
\mathbf F_2^{\,h-1}\rtimes\mathrm{GL}_{h-1}(\mathbf F_2).
\]
\end{theorem}

Together with the zero- and rank-one cases and the plane theorem, this gives every geometric and arithmetic image for projective hypersurfaces in the stated ordinary sense. The ordinary basis cover is geometrically connected in the full linear cases. In the other characteristic-two parity, the canonical covector gives the precise obstruction and its entire stabilizer occurs. The residue threshold retains the same value in both cases. The geometric image is computed after base change to an algebraic closure of the perfect ground field. The arithmetic image contains it and lies in the ambient general linear group; in the exceptional parity it also preserves the descended canonical covector. These inclusions force the stated arithmetic equalities.

\paragraph{Credit and context.}
The projective Hasse--Witt coefficient matrix is Katz's formula \citep[\S2.3.7, pp.~34--36, especially Algorithm~2.3.7.14 and formula~2.3.7.16]{Katz1972}. Katz also proves the general point-count congruence and the relation of stable coherent Frobenius to \'etale $\mathbf F_p$ cohomology \citep[Expos\'e XXII, Proposition~2.0 and Corollary~2.1, pp.~406--407; Theorem~3.1 and Corollary~3.2, pp.~410--411]{Katz}. Generic invertibility is classical: \citet[Theorem~3, p.~142; proof pp.~142--144]{Koblitz} proves it for general complete intersections of every multidegree and explicitly credits Miller's earlier hypersurface result. For the precise combinatorial form used here, \citet[Theorem~2.3, p.~2; Proposition~2.11, p.~3 of the arXiv version]{AdolphsonSperber} establish that an explicitly normalized integral determinant has constant term $1$. We use that theorem with its exact support hypothesis. Our coefficient calculation fixes the transpose and arithmetic/geometric Frobenius conventions. Projective cohomology and Artin--Schreier are the standard results in \citet[Lemma~30.8.1 and Section~59.63]{StacksProjective,StacksAS}. No priority claim is made for these classical inputs or for the formulations in this paper.

Huang, Lian, Yau, and Yu extend the coefficient and generic-invertibility formulas to toric hypersurfaces \citep[author version, Corollaries 2.3 and 2.8, pp.\ 8--11]{HLYY}. Their Theorem 5.3, pp.\ 19--22, reconstructs the unit-root Frobenius and connection matrices under the Hodge-freeness, Hasse--Witt and toric-vertex hypotheses stated there. This concerns the matrices and their $p$-adic lifts; the finite monodromy images in Theorems~\ref{thm:plane-images} and~\ref{thm:higher-images} require the geometric and group arguments below.

The recent preprint of \citet[Theorem~1.1, pp.~1--2]{OpenAI} states that, after good reduction of an abelian variety, pairing a specialized rational Hodge class with complementary products of divisor classes gives the same rational number in crystalline and every prime-to-$p$ realization. It provides neighboring context on compatible rational cohomological pairings; none of its results is used in this paper.

\paragraph{The proof mechanism.}
Write $d=m+3$ and specialize to a cubic together with $m$ hyperplanes. The top cohomology has a filtration by constant spaces, rank-two quotients of three-point permutation modules, and ordinary elliptic lines. Its rank is exactly the generic rank. A double-root trait gives a permutation supported on one node set, even in characteristic two; an isolated simple Hasse zero gives a scalar supported on one elliptic line. Their traces exhaust the prime field. The finite Lang system extends over a normal coefficient open containing this boundary, so the trace conclusion passes to the smooth generic family without a rank-drop extension. A finite-cover twist and Lang--Weil then produce the required points in every nonempty open. The characteristic-two matrix and covector calculations also support the exact plane groups. For those groups, a moving line or conic gives the sum-zero monodromy of lifted intersection points on an ordinary curve. Gluing their torsors produces actual transvections in the full cohomology extension. At a line arrangement these subgroups generate the required linear group or covector stabilizer. For $p\ge5$, the cubic scalar actions can be located in the alternating parts of an arrangement representation. Explicit overlap coefficients give full linear groups on those parts, and an elementary amalgamation argument gives the entire higher-dimensional group. In characteristic three at $d=n+2$, a coefficient partition gives a regular chart with the entire Hasse--Witt matrix as free coordinates. In larger degrees, the modular cubic permutation module forces supported transvections through a three-cycle norm. In characteristic two, an explicit elliptic Artin--Schreier cover and a selected intersection subcomplex realize the same kind of transformations on the full higher-dimensional extension. Their coordinate triple operations generate the precise linear group or covector stabilizer.

\section{Cohomology, coordinates, and the Hasse--Witt matrix}\label{sec:pn-1}

Katz gives the Hasse--Witt formula for smooth hypersurface families over an arbitrary base \citep[\S2.3.7, pp.~33--36, especially Algorithm~2.3.7.14]{Katz1972}. The short coordinate derivation below fixes our row and column convention and also applies to singular hypersurfaces. That larger form is needed for the cubic and hyperplane boundary.

Let \(n\geq2\), let \(d\geq1\), and let
\(Y=V(f)\subset\mathbf P^n_k\) be a degree-\(d\) hypersurface, with
\(f\neq0\). Write

\begin{equation}
 I_d=\{a=(a_0,\ldots,a_n)\in\mathbf Z_{>0}^{n+1}:|a|=d\},\qquad
 h=|I_d|=\binom{d-1}{n},                                                 \label{eq:pn-1-1}
\end{equation}

where \(h=0\) if \(d\leq n\). The standard affine Čech cover of projective
space gives

\begin{equation}
 H^n(\mathbf P^n,\mathcal O(-d))
 =\bigoplus_{a\in I_d} kx^{-a},
 \qquad H^i(\mathbf P^n,\mathcal O(j))=0\quad(0<i<n).                     \label{eq:pn-1-2}
\end{equation}

This computation, over an arbitrary coefficient ring as well as a field, is the projective-space calculation of \citet[Lemma 30.8.1, Tag 01XT]{StacksProjective}. One can see the displayed basis directly: the top
Čech quotient retains precisely Laurent monomials whose every exponent
is negative.

The sequence

\begin{equation}
 0\longrightarrow\mathcal O(-d)\xrightarrow{\ f\ }
 \mathcal O\longrightarrow\mathcal O_Y\longrightarrow0
\end{equation}

therefore identifies

\begin{equation}
 H^{n-1}(Y,\mathcal O_Y)\simeq H^n(\mathbf P^n,\mathcal O(-d)),            \label{eq:pn-1-3}
\end{equation}

and gives \(H^0(Y,\mathcal O_Y)=k\) and
\(H^i(Y,\mathcal O_Y)=0\) for \(0<i<n-1\).

In characteristic \(p\), the map on the left term of this sequence that
commutes with Frobenius on \(\mathcal O\) is
\(u\mapsto f^{p-1}u^p\). Thus, with columns and rows indexed by \(I_d\),
the semilinear Frobenius is \(v\mapsto B_p(f)v^{[p]}\), where

\begin{equation}
 B_p(f)_{b,a}=[x^{pa-b}]f^{p-1}.                                        \label{eq:pn-1-4}
\end{equation}

A coefficient with a negative exponent is zero. More generally, for
\(q=p^r\), the \(r\)-fold Frobenius has matrix

\begin{equation}
 B_q(f)_{b,a}=[x^{qa-b}]f^{q-1},\qquad
 B_q=B_pB_p^{[p]}\cdots B_p^{[p^{r-1}]}.                                 \label{eq:pn-1-5}
\end{equation}

The same argument on the universal projective-space Čech complex works
relatively over every coefficient chart. On rescaling the equation from
\(f\) to \(uf\), coherent coordinates change by \(v\mapsto u^{-1}v\), and
\(B_p\mapsto u^{p-1}B_p\). These transitions give the intrinsic universal
Frobenius map. In particular, invertibility of \(B_p\) is independent of
the equation used for the hypersurface.

\section{The Lang local system and Frobenius conventions}\label{sec:pn-2}

This section constructs its finite local system and fixes the Frobenius convention.

The Artin--Schreier sequence on the étale site is

\begin{equation}
 0\to\mathbf F_p\to\mathcal O\xrightarrow{F-1}\mathcal O\to0.             \label{eq:pn-2-1}
\end{equation}

Étale and Zariski cohomology agree for the structure sheaf. These are the standard inputs in \citet[Section 59.63, Tag 0A3J]{StacksAS}.
The preceding cohomology calculation, together with surjectivity of
\(x\mapsto x^p-x\) over an algebraically closed field when \(n=2\), gives

\begin{equation}
 H^{n-1}_{\mathrm{et}}(Y,\mathbf F_p)
 =\ker\bigl(F-1:H^{n-1}(Y,\mathcal O_Y)\to H^{n-1}(Y,\mathcal O_Y)\bigr).  \label{eq:pn-2-2}
\end{equation}

For clarity, the elementary fixed-vector fact needed below is recorded in
the stronger prime-power version. If \(Q\) is any power of \(p\) and
\(B\in\mathrm{GL}_h(\bar k)\), the equations

\begin{equation}
                       Bv^{[Q]}=v                                    \label{eq:pn-2-3}
\end{equation}

cut out a finite étale group with \(Q^h\) geometric points. Indeed,
multiplying by \(B^{-1}\) gives equations
\(v_i^Q-\sum_j(B^{-1})_{ij}v_j=0\). Their leading monomials for a degree
order are the pairwise relatively prime \(v_i^Q\); division leaves the
monomials \(v_1^{e_1}\cdots v_h^{e_h}\), \(0\leq e_i<Q\), as a basis.
The relations' pairwise reductions agree because products of the linear
right sides commute. The algebra is therefore free of rank \(Q^h\).
Its Jacobian is \(-B^{-1}\), so it is étale.

The fixed points form an \(\mathbf F_Q\)-vector space of dimension \(h\)
and span \(\bar k^h\). To check the last assertion, choose a maximal
\(\bar k\)-independent list of fixed vectors, of length \(s\). In its
span, a vector is fixed exactly when its coefficients are in
\(\mathbf F_Q\); hence that span contains only \(Q^s\) fixed vectors.
All fixed vectors lie in it by maximality, so \(Q^s=Q^h\), and \(s=h\).

With \(Q=p\), this construction over a base on which \(B_p\) is invertible
gives the rank-\(h\) finite étale \(\mathbf F_p\)-local system in \eqref{eq:pn-2-2}.
It also gives its finite étale basis torsor. More explicitly, the vector
scheme is finite locally free by the monic calculation above. In its
\(h\)-fold product, the basis locus is open and closed: if
\(T=B_pT^{[p]}\), \(t=\det T\), and \(b=\det B_p\), then
\(t=bt^p\), and \(e=bt^{p-1}\) is idempotent. The locus \(t\neq0\)
is exactly \(e=1\). The basis torsor has group
\(\mathrm{GL}_h(\mathbf F_p)\). Relative Artin--Schreier and the same
relative Čech computation identify it with the hypersurface local system;
no prime-to-\(p\) smooth-base-change assertion is used.

If \(f\) is defined over \(\mathbf F_q\), let \(T\) be a matrix of fixed
basis vectors. Arithmetic Frobenius on this basis is the matrix \(g\)
specified by

\begin{equation}
 T^{[q]}=Tg,\qquad g\in\mathrm{GL}_h(\mathbf F_p).
\end{equation}

Since \(B_qT^{[q]}=T\), it follows that

\begin{equation}
                         B_q=Tg^{-1}T^{-1}.                            \label{eq:pn-2-4}
\end{equation}

Thus the trace in \eqref{eq:pn-1-5} is the trace of \textbf{geometric} Frobenius on the
\(\mathbf F_p\)-local system. Confusing it with arithmetic Frobenius would
invert the matrix. Inversion permutes a finite monodromy group, but the
convention matters for the stated point-count identity.

\section{The known generic-invertibility theorem}\label{sec:pn-3}
The nonemptiness of the ordinary open follows from a theorem of Adolphson--Sperber. We state the precise form used, including its support condition.

Let $f=\sum_e u_e x^e$ be a degree-$d$ polynomial with indeterminate coefficients, and suppose its support contains every exponent in $I_d$. Adolphson--Sperber form the integral matrix
\[
A_{a,b}=[x^{pa-b}]f^{p-1}\qquad(a,b\in I_d).
\]
This is the transpose of our matrix $B_p$. Normalize its rows and columns by
\[
\widetilde A_{a,b}=u_a^{-p}u_b A_{a,b}.
\]
Their Proposition~2.11 proves that $\det\widetilde A$ has constant term exactly $1$, both integrally and after reduction modulo $p$. Equivalently, the coefficient of $\prod_{a\in I_d}u_a^{p-1}$ in $\det B_p$ is exactly $1$. Their Theorem~2.3 therefore gives
\[
\det B_p\ne0\quad\hbox{in the universal coefficient ring over }\mathbf F_p.
\]
The support hypothesis holds for the universal degree-$d$ polynomial. The exact statements and full combinatorial proof are in \citet[Section~2, pp.~2--4 of the arXiv version]{AdolphsonSperber}. For $h=0$, the determinant is $1$ by convention.

The smooth hypersurface open is nonempty in every characteristic. At a fixed point of $\mathbf P^n$, degree-$d$ forms span the value and the $n$ first local derivatives: on a coordinate chart, their first jets include a constant and every linear coordinate. Singularity at that point therefore imposes $n+1$ independent linear conditions. The incidence of a point and a hypersurface singular there has dimension at most one less than the hypersurface parameter space. Its projection is closed because $\mathbf P^n$ is proper. Thus its complement is nonempty. In the integral parameter space, this smooth open intersects the nonvanishing open of $\det B_p$. This proves the nonemptiness and density used throughout. It gives no full monodromy conclusion by itself.

\section{Exact point-count congruence and the lower threshold}\label{sec:pn-4}

This section counts points directly and proves the obstruction below the sharp threshold.

Let \(q=p^r\) and \(f\in\mathbf F_q[x_0,\ldots,x_n]\) be homogeneous of
degree \(d\). Write \(N=\#V(f)(\mathbf F_q)\) in projective space and
\(Z\) for the number of affine zeros in \(\mathbf F_q^{n+1}\). Then

\begin{equation}
                     Z=1+(q-1)N.                                     \label{eq:pn-4-1}
\end{equation}

Working in \(\mathbf F_q\), the characteristic function of the affine
zero set is \(1-f^{q-1}\), and hence
\(Z=-\sum_{x\in\mathbf F_q^{n+1}}f(x)^{q-1}\).
For one variable, the sum of its \(e\)-th powers is zero if \(e=0\),
zero if \(e>0\) is not divisible by \(q-1\), and \(-1\) if \(e>0\)
is divisible by \(q-1\). Therefore only exponents
\((q-1)a\), \(a\in I_d\), survive. Combining this with \eqref{eq:pn-4-1} gives

\begin{equation}
 \boxed{\quad
 \#Y(\mathbf F_q)\equiv
 1+(-1)^{n-1}\sum_{a\in I_d}[x^{(q-1)a}]f^{q-1}
 =1+(-1)^{n-1}\operatorname{tr}B_q(f)\pmod p.
 \quad}                                                               \label{eq:pn-4-2}
\end{equation}

This calculation does not require smoothness or ordinarity.

If \(d\leq n\), the index set \(I_d\) is empty, so every hypersurface
has residue 1 over every finite extension. Its top coherent cohomology
vanishes, and \eqref{eq:pn-2-2} also proves the zero-rank/trivial-monodromy assertion.

If \(d=n+1\), there is a single index \(a=(1,\ldots,1)\). Put

\begin{equation}
 A(f)=[(x_0\cdots x_n)^{p-1}]f^{p-1}.
\end{equation}

On the nonempty ordinary open, \(A(f)\neq0\). Over \(\mathbf F_q\),
the scalar \(B_q\) is
\(A(f)A(f)^p\cdots A(f)^{p^{r-1}}\), which is nonzero.
Equation \eqref{eq:pn-4-2} consequently excludes residue 1 on this entire open.
This proves both lower-threshold statements without using a full-image assertion.

\section{Rank-one geometric monodromy in odd characteristic}\label{sec:pn-5}

This section determines the full geometric image in the one-dimensional case.

Assume \(d=n+1\) and \(p>2\), and put \(m=(p-1)/2\). On an affine
coefficient chart, the basis torsor for the rank-one local system is

\begin{equation}
                   z^{p-1}=A(f)^{-1}.                                 \label{eq:pn-5-1}
\end{equation}

It suffices to show that the rational function \(A\) has a divisorial
zero of order 1 over the algebraic closure of the ground field. Then
\(y^{p-1}-A\) is Eisenstein at that divisor, so \eqref{eq:pn-5-1}, after replacing
\(z\) by \(1/y\), has geometric degree \(p-1\).

Consider the one-parameter forms

\begin{equation}
 f_\lambda=
 \bigl(Y^2Z-X(X-Z)(X-\lambda Z)\bigr)x_3\cdots x_n,                       \label{eq:pn-5-2}
\end{equation}

where the product is 1 if \(n=2\). They lie in the affine coefficient
chart in which the coefficient of \(Y^2Zx_3\cdots x_n\) is 1.
Choosing exactly \(m\) copies of \(Y^2Z\) in the \(2m\)-th power gives

\begin{equation}
 A(f_\lambda)=\binom{2m}{m}H_m(\lambda),\qquad
 H_m(\lambda)=\sum_{i=0}^m\binom mi^2\lambda^i.                          \label{eq:pn-5-3}
\end{equation}

For the coefficient calculation, after removing \(Y^{2m}Z^mX^m\), one
needs the coefficient of \(X^mZ^m\) in
\((X-Z)^m(X-\lambda Z)^m\). The two signs cancel, and symmetry under
\(i\leftrightarrow m-i\) gives \eqref{eq:pn-5-3}.

The polynomial \(H_m\) has degree \(m\geq1\),
\(H_m(0)=1\), and
\(H_m(1)=\binom{2m}{m}\equiv(-1)^m\neq0\pmod p\).
Its coefficient recurrence

\begin{equation}
 (i+1)^2\binom m{i+1}^2=(m-i)^2\binom mi^2
\end{equation}

implies, in characteristic \(p\),

\begin{equation}
 \lambda(1-\lambda)H_m''+(1-2\lambda)H_m'-m^2H_m=0.                     \label{eq:pn-5-4}
\end{equation}

Indeed \(2m=-1\) in the field, so the coefficient recurrence obtained
from \eqref{eq:pn-5-4} is exactly the displayed one. A repeated root \(\lambda_0\)
cannot be 0 or 1. If its multiplicity is \(j\), then
\(2\leq j\leq m<p\); the lowest-order term of the first summand of
\eqref{eq:pn-5-4} has nonzero coefficient
\(\lambda_0(1-\lambda_0)j(j-1)\), while the other terms have higher
order. This is impossible. Thus every root of \(H_m\) is simple.

At a root, the derivative of the restriction of the universal polynomial
\(A\) to \eqref{eq:pn-5-2} is nonzero. Hence its zero divisor on that smooth affine
coefficient space has a component of multiplicity 1. The forms in \eqref{eq:pn-5-2}
need not themselves be smooth when \(n>2\): the argument uses a valuation
of the ambient function field, and does not require the divisor to lie
inside the smooth ordinary open.

The Eisenstein argument proves geometric degree \(p-1\) of the basis
cover at the generic point. Restriction to the nonempty smooth ordinary
open preserves this degree and connectedness. Its geometric monodromy
is therefore \(\mathbf F_p^*=\mathrm{GL}_1(\mathbf F_p)\). The arithmetic
image contains it and is already a subgroup of the same group, so it is
equal to it over every perfect ground field. In characteristic two, the
rank-one target group \(\mathrm{GL}_1(\mathbf F_2)\) is trivial.

\section{Linear restrictions and plane-cubic degenerations}\label{sec:cubic-geometry}
These extension and plane-cubic calculations provide the simultaneous degenerations used to isolate individual monodromy factors.


\subsection{Compatible restrictions on two linear subspaces}\label{sec:cubic-restrictions}

If \(A,B\subset\mathbf P^n\) are linear subspaces, every pair of cubic forms on \(A\) and \(B\) agreeing on \(A\cap B\) extends to an ambient cubic. Choose coordinates \(z\) on the vector-space intersection, coordinates \(x\) completing it in \(A\), and coordinates \(y\) completing it in \(B\). If the forms are \(f(x,z)\) and \(g(y,z)\), with common restriction \(h(z)\), then

\[
f(x,z)+g(y,z)-h(z)
\]

extends the pair to the span of \(A\cup B\), and then to the ambient space. For disjoint projective subspaces the vector-space intersection is zero, and the same formula is \(f+g\). The same argument applies when one subspace contains the other.

Consequently, restriction of projective ambient cubic forms to a fixed linear subspace, away from zero restriction, is an affine-space bundle over the projective space of forms on that subspace. To see the bundle statement, locally normalize a nonzero coefficient of the restricted form to 1 and choose a linear splitting of the restriction map. The remaining coordinates are precisely its kernel. In particular, the inverse image of an irreducible locus of nonzero restricted forms is irreducible.

\subsection{Smooth extensions with a prescribed restriction}\label{sec:cubic-extensions}

Let \(A\) be a proper linear subspace of \(W=\mathbf P^r\), and prescribe a cubic form on \(A\). The extensions form a nonempty affine space. At \(x\in W\setminus A\), choose a linear form \(l\) vanishing on \(A\) with \(l(x)\ne0\). The variations \(l^2L\), as \(L\) ranges over all linear forms, belong to the kernel of restriction and span all first jets at \(x\). Indeed multiplication by the unit \(l^2\) is an isomorphism on the first-jet space, and linear forms span that space. This is valid also in characteristic 2.

The condition that an extension and all its local first derivatives vanish at a fixed \(x\notin A\) therefore imposes \(r+1\) independent affine-linear conditions. The incidence of such pairs (extension, singular point) has dimension at most

\[
\dim(\text{extension space})-(r+1)+r
=\dim(\text{extension space})-1.
\]

The closure of its projection is a proper closed subset. A generic extension is thus smooth away from \(A\).

At a zero of the prescribed restriction where that restriction cuts out a smooth hypersurface of \(A\), a derivative tangent to \(A\) is nonzero, so every extension is smooth there. If \(A\) is a point on which the prescribed value is nonzero, there are no zeros on \(A\). Finally, if \(A\) is a line whose prescribed binary cubic has one double zero and one distinct simple zero, the extension is automatically smooth at the simple zero. At the double zero, for \(r\ge2\), the normal first derivatives can be chosen freely using forms in the ideal of \(A\). Their simultaneous vanishing is a proper linear condition. Thus a generic extension is smooth there as well.

These arguments also show that unrestricted cubic hypersurfaces in every positive-dimensional projective space have a nonempty smooth locus: with no prescribed subspace, the full cubic system spans every first jet. For \(\mathbf P^1\) this is the existence of a squarefree binary cubic. In dimension zero the required transversality is simply a nonzero value.

All these smoothness and transversality conditions are open. The bad incidence can equivalently be formed over the projective space \(W\), so its projection is closed. We shall repeatedly use the following elementary consequence: on one irreducible family, finitely many open conditions hold simultaneously if each has a nonempty locus. The examples establishing the separate conditions need not be the same member of the family.

\subsection{Plane Hasse invariants}\label{sec:plane-cubic-hasse}

The coefficient formula \eqref{eq:pn-1-4}, with \(n=2,d=3\), identifies the plane Hasse invariant with the coefficient of \((XYZ)^{p-1}\) in \(f^{p-1}\). In odd characteristic, the Legendre calculation of Section~\ref{sec:pn-5} gives both smooth ordinary cubics and smooth Hasse-zero cubics. Its simple-root proof shows that the plane Hasse divisor has a smooth reduced point within the smooth-cubic locus: the derivative along the Legendre family is nonzero there. The rank-one basis cover is \eqref{eq:pn-5-1}.

For \(p=2\), use

\begin{equation}
f=Y^2Z+XYZ+X^3+Z^3.
\label{eq:cubic-3-3}
\end{equation}

The Hasse coefficient is 1. Its partial derivatives are

\[
f_X=YZ+X^2,\quad f_Y=XZ,\quad
f_Z=Y^2+XY+Z^2.
\]

If a singular point had \(Z=0\), the first derivative would force \(X=0\), and then \(f_Z=Y^2\ne0\). If \(Z\ne0\), the equation \(f_Y=0\) forces \(X=0\), then \(f_X=0\) forces \(Y=0\), contradicting \(f=0\). Thus \eqref{eq:cubic-3-3} is smooth ordinary.

\subsection{Ordinary cubics with a simple tangent}\label{sec:plane-cubic-tangent}

In characteristic 2, \eqref{eq:cubic-3-3} restricts to the line \(X=0\) as

\[
Z(Y+Z)^2.
\]

It therefore has one double intersection point and a distinct simple intersection point on that line.

For odd \(p\), choose an ordinary Legendre cubic and write its affine equation \(y^2=g(x)\). At a point with \(y\ne0\), the tangent slope is \(s=g'(x)/(2y)\). The quadratic coefficient after restricting to this tangent is \(s^2-g''(x)/2\). Its numerator, after multiplication by \(4g(x)\), is

\begin{equation}
g'(x)^2-2g(x)g''(x).
\label{eq:cubic-3-4}
\end{equation}

For \(p\ne3\) this polynomial has leading coefficient \(-3\), so it is nonzero. For \(p=3\), its cubic coefficient is \(-2(1+\lambda)\), which is nonzero for an ordinary Legendre cubic because its Hasse polynomial is \(1+\lambda\). Choose \(x\) avoiding the finitely many zeros of \(g\) and \eqref{eq:cubic-3-4}. The tangent has intersection multiplicity exactly 2 at the selected point, and its third intersection is distinct. Projective changes of coordinates move this tangent to any prescribed line.

\subsection{Prescribed squarefree binary cubics}\label{sec:plane-cubic-restriction}

Every smooth plane cubic has a transverse line. Indeed its tangent lines form the image of its one-dimensional Gauss map in the two-dimensional dual plane. Choose a line outside that image. A projective change of coordinates can move this line to any prescribed line.

Given any nonzero squarefree binary cubic on a fixed line, there is a smooth ordinary plane cubic with exactly that restriction. Start with one of the ordinary cubics above and a transverse line. The group \(\operatorname{PGL}_2\) on the line takes its three distinct intersection points to any prescribed three distinct points. Extend that line automorphism to a projective transformation of the plane and then rescale the equation to match the prescribed binary form. Likewise an ordinary plane cubic can be chosen with any prescribed nonzero value at a fixed point: first move it to avoid the point and then rescale.

These assertions concern exact restricted forms, not merely their zero sets; the final scalar adjustment is what gives exact compatibility in Section~\ref{sec:cubic-restrictions}.

\subsection{A component of the Hasse-zero divisor}\label{sec:plane-cubic-component}

For odd \(p\), take the irreducible component \(D\) of the plane-cubic Hasse-zero locus passing through a smooth reduced point furnished by Section~\ref{sec:plane-cubic-hasse}. Work throughout in the open of smooth plane cubics. The component \(D\) is irreducible, and its smooth reduced locus is nonempty open.

The connected group \(\operatorname{PGL}_3\) preserves \(D\). Here is an explicit component argument: the closure of the image of \(\operatorname{PGL}_3\times D\) is irreducible, is contained in the Hasse-zero locus, and contains the irreducible component \(D\). Maximality of that component forces equality. Therefore we can move a fixed smooth Hasse-zero cubic within \(D\) to meet any given line transversely or to avoid any given point. Moving a smooth reduced point also preserves the smooth reduced condition.

\section{Independently supported cubic-boundary monodromy}\label{sec:cubic-joint}
To prove Theorem~\ref{thm:joint}, we construct monodromy supported on each factor while the other factors extend over the same trait.


\subsection{Nonemptiness of the common open}\label{sec:cubic-open}

Every individual stratum-smoothness condition is nonempty by Section~\ref{sec:cubic-extensions}. Every individual plane-ordinarity condition is nonempty because restriction to that plane is surjective and ordinary plane cubics exist by Section~\ref{sec:plane-cubic-hasse}. Avoidance of each zero-dimensional stratum is also a nonempty open condition. Their finite intersection in the irreducible ambient cubic parameter space is nonempty. This proves \(S\ne\varnothing\).

\subsection{An isolated double-root degeneration}\label{sec:cubic-node-degeneration}

Fix one \(J\), and abbreviate \(\ell_J\) to \(\ell\). In the projective space of nonzero binary cubics, the closure of the locus

\[
L_a^2L_b,\qquad a\ne b\in\mathbf P^1,
\]

is an irreducible divisor. This remains true in characteristic 2: it is the two-dimensional image of the irreducible surface \(\mathbf P^1\times\mathbf P^1\); its general member has one double root and one distinct simple root. No assertion about the reducedness of a discriminant polynomial is needed. Let \(D_J\) be its inverse image under restriction of ambient cubics to \(\ell\), omitting zero restriction. By Section~\ref{sec:cubic-restrictions}, \(D_J\) is irreducible.

There is a nonempty open subset of \(D_J\) on which all of the following hold: the restriction to \(\ell\) has exactly one double root and one distinct simple root, both away from the fixed zero-dimensional strata; every other \(Z_{J'}\) is reduced; all the \(E_I\) are smooth ordinary; and \(Q\) is smooth and transverse to every other hyperplane stratum. We verify that each requested condition is nonempty on this one irreducible family.

For another line \(\ell_{J'}\), the intersection with \(\ell\) is empty or a point. Choose the double-root form on \(\ell\) to be nonzero at that point, if present. Choose a squarefree binary cubic on \(\ell_{J'}\) with the same nonzero value there, and use Section~\ref{sec:cubic-restrictions} to extend the pair. This proves nonemptiness of the reducedness condition for \(Z_{J'}\).

For a plane \(\Lambda_I\) containing \(\ell\), Section~\ref{sec:plane-cubic-tangent} supplies a smooth ordinary cubic with the required double-plus-simple line section. If the plane meets \(\ell\) in a point, choose an ordinary plane cubic avoiding that point and a double-plus-simple binary cubic with the same nonzero value. If they are disjoint, choose both freely. Section~\ref{sec:cubic-restrictions} extends the compatible pair. Thus each plane-ordinarity condition has a nonempty locus on \(D_J\).

For any other positive-dimensional hyperplane stratum \(W\), its intersection with \(\ell\) is \(\ell\), a point, or empty. If it contains \(\ell\) and has dimension at least 2, the double-root case of Section~\ref{sec:cubic-extensions} gives a smooth extension to \(W\). If the intersection is a point, choose a nonzero restricted value and use the point case. If empty, use unrestricted smoothness. Extension from \(W\cup\ell\) to the ambient space then follows from Section~\ref{sec:cubic-restrictions}. The stratum \(W=\mathbf P^n\) is included. Zero-dimensional strata are avoided by a nonempty open condition. This verifies every condition separately and hence simultaneously.

Choose a member \(Q_0\) of this common open. Use coordinates on \(\ell\) so that its restricted cubic is \(X^2(X-Z)\), after scaling. Let \(R=\kappa[[t]]\), \(K=\kappa((t))\). Lift the following deformations of its binary restriction to ambient cubics using a linear splitting of the restriction map:

\begin{equation}
\begin{array}{ll}
p\ne2:&(X^2-tZ^2)(X-Z),\\[2mm]
p=2:&(X^2+tXZ+tZ^2)(X-Z).
\end{array}
\label{eq:cubic-4-1}
\end{equation}

The lifted family has special member \(Q_0\), and its generic member lies in \(S\). Every other condition already holds at the special member and remains true over the trait.

For odd \(p\), \(x^2-t\) defines a separable quadratic extension of \(K\). For \(p=2\), \(x^2+tx+t\) is Eisenstein and separable; its roots are \(\alpha\) and \(\alpha+t\), so its degree-two extension is Galois. The third root 1 is fixed and distinct. Hence local inertia gives a transposition of the selected three-point set.

All other node covers and all elliptic rank-one local systems extend as finite étale objects over \(R\): their special fibres are respectively reduced or smooth ordinary. Since \(R\) is strictly henselian, their local monodromy is trivial. The resulting element of the joint image is therefore a transposition supported only on \(Z_J\). This argument includes the wild quadratic in characteristic 2.

\subsection{From one transposition to a supported symmetric group}\label{sec:cubic-node-group}

The projection of the joint image to each node factor is \(S_3\). To prove this without assuming generic polynomial Galois theory, consider the ordered-root cover of the squarefree binary-cubic locus. Its total space is

\[
\operatorname{Conf}_3(\mathbf P^1)
=\{(a,b,c):a,b,c\text{ pairwise distinct}\},
\]

which is irreducible. The map forgetting the ordering is a finite étale \(S_3\)-torsor in every characteristic: the three roots are distinct, and the six possible orderings give the étale splitting cover. Pulling back by ambient restriction gives an affine-space bundle over this irreducible cover. Its restriction over the nonempty open \(S\) stays nonempty and irreducible. Thus its monodromy is all \(S_3\).

Conjugate the supported transposition from Section~\ref{sec:cubic-node-degeneration} by joint-image elements lifting arbitrary permutations in this \(S_3\). The other coordinates of each conjugate remain the identity. All transpositions in the target factor are obtained, and they generate \(S_3\). This proves the required supported node factor.

\subsection{An isolated elliptic Hasse-zero degeneration}\label{sec:cubic-scalar}

Assume \(p\) odd and fix a target plane \(\Lambda_I=\Lambda\). Pull back the component \(D\) of Section~\ref{sec:plane-cubic-component} under restriction to \(\Lambda\). Its inverse image \(T_I\) is irreducible by Section~\ref{sec:cubic-restrictions}.

There is a member \(Q_0\) of \(T_I\) for which the target cubic is a smooth reduced point of its Hasse-zero divisor, every other plane cubic is smooth ordinary, all node sets are reduced, and all hyperplane-stratum transversality conditions hold. Again we verify the separate nonempty open conditions on this one irreducible family.

For another plane \(\Lambda_{I'}\), its intersection with \(\Lambda\) is a line, a point, or empty. In the line case, move a smooth reduced member of \(D\) to make its line restriction squarefree, using Section~\ref{sec:plane-cubic-component}. Section~\ref{sec:plane-cubic-restriction} gives an ordinary cubic on the other plane with exactly that binary restriction. Extend the pair by Section~\ref{sec:cubic-restrictions}. In the point case, move the target cubic to avoid the point and choose an ordinary cubic on the other plane with the same nonzero value. The empty-intersection case imposes no compatibility. Therefore the ordinary locus for every other plane is a nonempty open in \(T_I\).

For any linear stratum \(W\), its intersection with \(\Lambda\) is \(\Lambda\), a line, a point, or empty. If \(W\) contains \(\Lambda\), a smooth target cubic extends smoothly to \(W\) by Section~\ref{sec:cubic-extensions}. If the intersection is a line, move the target so that its restriction is squarefree and apply that lemma. If the intersection is a point, move the target to avoid it. If \(W\) is contained in \(\Lambda\), the same line or point argument applies directly. The empty case is unrestricted. These observations give every stratum condition, including all node sets and ambient smoothness, on a nonempty open. Finally, the reduced smooth part of the target Hasse divisor is itself nonempty open. Irreducibility proves the simultaneous assertion.

Take a formal trait through this \(Q_0\), transverse to the target Hasse divisor. A local trivialization of the Hasse line gives a scalar of valuation 1, say \(t\) times a unit. By the rank-one basis equation \eqref{eq:pn-5-1}, the target ordinary basis cover on the punctured trait is

\begin{equation}
z^{p-1}=\operatorname{Hasse}^{-1}.
\label{eq:cubic-4-2}
\end{equation}

Every unit has a \((p-1)\)-st root in \(\kappa[[t]]\), while the valuation of the right side is \(-1\). Thus \eqref{eq:cubic-4-2} has full cyclic inertia of order \(p-1\), identified with \(\mathbf F_p^\times\). Every other listed factor extends finite étale over the strictly henselian trait and has trivial inertia. We have obtained the entire scalar group supported only on the target elliptic factor.

For \(p=2\), each elliptic factor is a one-dimensional \(\mathbf F_2\)-space and its automorphism group is already trivial. The supported groups constructed above generate the full direct product of all listed factor groups. This proves Theorem~\ref{thm:joint}.

\section{The exact \v Cech filtration on the cubic boundary}\label{sec:cubic-cech}
The cohomology of the cubic and hyperplane union has an explicit Frobenius-stable filtration whose rank agrees with the smooth generic rank.


Set \(d=m+3\), and for \(Q\in S\) define

\begin{equation}
Y_0=Q\cup H_1\cup\cdots\cup H_m\subset\mathbf P^n.
\label{eq:cubic-5-1}
\end{equation}

It is a simple normal crossings hypersurface of degree \(d\). The hypotheses defining \(S\) ensure that the hyperplanes together with \(Q\) have independent local equations wherever they meet. In particular no \(n+1\) of these components meet.

\subsection{Exactness of the intersection complex}\label{sec:cubic-cech-resolution}

The alternating restriction complex

\begin{equation}
0\longrightarrow\mathcal O_{Y_0}\longrightarrow
\bigoplus_a\mathcal O_{D_a}\longrightarrow
\bigoplus_{a<b}\mathcal O_{D_a\cap D_b}\longrightarrow\cdots
\label{eq:cubic-5-2}
\end{equation}

is exact, where \(D_0=Q\) and \(D_i=H_i\). This can be checked in normal crossings coordinates. Locally the union is \(t_1\cdots t_r=0\). For a fixed monomial, let \(A\) be the nonempty set of indices whose exponent is zero; that monomial survives exactly on intersections indexed by subsets of \(A\). Its augmented coefficient complex is the simplex complex on \(A\), which is contractible. This proves exactness over a polynomial ring, and flat étale change of coordinates gives the local assertion on \(Y_0\). The same proof works with power series. Signs cause no difficulty in characteristic 2.

The hypercohomology spectral sequence of \eqref{eq:cubic-5-2} is compatible with Frobenius, because restriction and the alternating differentials commute with \(p\)-th powers.

\subsection{The nonzero cohomology rows}\label{sec:cubic-cech-rows}

A hyperplane-only intersection is a projective space and has no positive coherent cohomology. An intersection containing \(Q\) and \(a\) hyperplanes is a smooth cubic hypersurface in \(\mathbf P^{n-a}\). The projective cohomology calculation of Section~\ref{sec:pn-1} and the hypersurface exact sequence show:

\begin{itemize}
\tightlist
\item
  for \(a<n-2\), its positive coherent cohomology vanishes and \(H^0=\kappa\);
\item
  for \(a=n-2\), it is \(E_I\), with \(H^0=\kappa\) and \(H^1\) of dimension 1;
\item
  for \(a=n-1\), it is \(Z_J\), with \(H^0=\kappa^{Z_J}\) of dimension 3;
\item
  for \(a=n\), it is empty, by avoidance of the fixed points.
\end{itemize}

Thus the only positive-cohomology slot in the spectral sequence is

\[
E_1^{n-2,1}=\bigoplus_{|I|=n-2}H^1(\mathcal O_{E_I}).
\]

There is no incoming differential to this slot. Its only possible higher outgoing differential would land in Čech degree \(n\), but \eqref{eq:cubic-5-2} stops in degree \(n-1\). The zero-th cohomology row, denoted \(C^\bullet\), therefore gives a Frobenius-compatible exact sequence

\begin{equation}
0\longrightarrow H^{n-1}(C^\bullet)
\longrightarrow H^{n-1}(\mathcal O_{Y_0})
\longrightarrow\bigoplus_I H^1(\mathcal O_{E_I})
\longrightarrow0.
\label{eq:cubic-5-3}
\end{equation}

For \(n=2\), this is the same argument with the elliptic slot \((0,1)\) and the last Čech degree 1; no exceptional differential occurs.

\subsection{The constant row and node quotients}\label{sec:cubic-cech-quotients}

Every intersection of at most \(n\) components is nonempty. Replace \(H^0\) of each such intersection by its diagonal constant line. The resulting subcomplex \(K^\bullet\subset C^\bullet\) is the ordinary cochain complex of the full \((n-1)\)-skeleton of the simplex with \(m+1\) vertices. The quotient is concentrated in degree \(n-1\), where it is

\[
\bigoplus_{|J|=n-1}V_J,
\qquad V_J=\kappa^{Z_J}/\kappa\cdot\mathbf1.
\]

The skeleton is connected because \(n\ge2\), and its only reduced cohomology is in degree \(n-1\), of dimension \(\binom mn\). One direct proof is to augment the complex and compare with the full contractible simplex: it remains exact below the truncated top degree; the alternating sum of its term dimensions is

\[
(-1)^{n-1}\left(-1+\sum_{a=0}^{n-1}
                      (-1)^a\binom{m+1}{a+1}\right)
=\binom mn.
\]

The usual convention \(\binom mn=0\) for \(m<n\) includes the case of a full simplex rather than a proper skeleton.

The long exact sequence of the short exact sequence of complexes gives

\begin{equation}
0\longrightarrow\kappa^{\binom mn}
\longrightarrow H^{n-1}(C^\bullet)
\longrightarrow\bigoplus_J V_J\longrightarrow0.
\label{eq:cubic-5-4}
\end{equation}

The first term is canonically the base change of a constant \(\mathbf F_p\)-cochain cohomology space. Equations \eqref{eq:cubic-5-3} and \eqref{eq:cubic-5-4} are filtrations; no equivariant splitting is asserted or needed.

\subsection{Rank and invertibility}\label{sec:cubic-cech-rank}

The three types of graded pieces have ranks

\[
\binom mn,\qquad 2\binom m{n-1},\qquad\binom m{n-2}.
\]

Their sum is

\begin{equation}
\binom mn+2\binom m{n-1}+\binom m{n-2}
=\binom{m+2}n=\binom{d-1}n.
\label{eq:cubic-5-5}
\end{equation}

Frobenius is invertible on the constant part, on the permutation bundles of the finite étale node sets and their diagonal quotients, and on the elliptic pieces by ordinarity. Hence it is invertible on \(H^{n-1}(\mathcal O_{Y_0})\). All statements are compatible with families over \(S\): the stratum cohomology bundles have constant ranks, the displayed spectral sequence has just the stated rows, and the displayed exact sequences are exact sequences of locally free bundles with Frobenius.

\subsection{Fixed vectors and the boundary trace formula}\label{sec:cubic-lang-trace}

The relative Lang calculation of Section~\ref{sec:pn-2} shows that a rank-\(r\) bundle with invertible Frobenius has a finite étale fixed-vector local system of rank \(r\).

Taking fixed-vector sheaves is exact on exact sequences of these invertible Frobenius bundles. For surjectivity, lift a fixed section of a quotient locally to \(b\); then \(F(b)-b\) lies in the subbundle. Solve \(F(a)-a=F(b)-b\) étale-locally in that subbundle, and \(b-a\) is a fixed lift.

Apply this to \eqref{eq:cubic-5-3} and \eqref{eq:cubic-5-4}. The fixed vectors in the node bundle are the permutation module \(\mathbf F_p^{Z_J}\), and its quotient contributes \(\mathbf F_p^{Z_J}/\mathbf F_p\mathbf1\). The elliptic fixed-vector lines are naturally \(H^1_{\mathrm{et}}(E_I,\mathbf F_p)\). The latter identification follows from the Artin--Schreier sequence and coherent/étale comparison \citep[Section~59.63, Theorem~59.22.4, and Lemma~59.63.6]{StacksAS,StacksComparison,StacksProperAS}; the same identification is the rank-one basis-cover calculation in Section~\ref{sec:pn-5}.

For \(\gamma\in\pi_1(S)\), write \(a_I(\gamma)\in\mathbf F_p^\times\) for its action on the elliptic factor. Additivity of trace on an invariant filtration gives the exact formula

\begin{equation}
\operatorname{Tr}(\gamma\mid H^{n-1}(\mathcal O_{Y_0})^{F=1})
=\binom mn+
\sum_{|J|=n-1}\bigl(\#\operatorname{Fix}(\gamma\mid Z_J)-1\bigr)
+\sum_{|I|=n-2}a_I(\gamma)
\quad\text{in }\mathbf F_p.
\label{eq:cubic-6-1}
\end{equation}

In characteristic 3 the diagonal line of a three-point permutation module may lie in its augmentation submodule, but it still has dimension 1 and trace 1. Thus the quotient trace is still the number of fixed points minus 1. No division by 3, semisimplicity, or splitting occurs anywhere in \eqref{eq:cubic-6-1}.

\section{The generic local system and all traces}\label{sec:cubic-generic}

Here the comparison can be made directly on an open of the full parameter space; no entire generic local system is extended across a rank-drop locus.

Let \(\mathcal P_d=S_d=\mathbf P H^0(\mathbf P^n,\mathcal O(d))\), and let \(f:\mathcal Y\to\mathcal P_d\) be its universal hypersurface family. The universal exact sequence is

\[
0\to\mathcal O_{\mathbf P^n}(-d)\boxtimes
                 \mathcal O_{\mathcal P_d}(-1)
\to\mathcal O_{\mathbf P^n\times\mathcal P_d}
\to\mathcal O_{\mathcal Y}\to0.
\]

The projective-space cohomology computation of Section~\ref{sec:pn-1} gives a locally free bundle

\[
\mathcal M=R^{n-1}f_*\mathcal O_{\mathcal Y}
\cong H^n(\mathbf P^n,\mathcal O(-d))
             \otimes\mathcal O_{\mathcal P_d}(-1)
\]

of rank \(r=\binom{d-1}n\), valid over the full parameter space, including reducible hypersurfaces. It commutes with base change by this explicit calculation. Relative Frobenius is an algebraic Frobenius-linear map on \(\mathcal M\). Let \(W\subset\mathcal P_d\) be its invertibility locus. This is a smooth normal integral open subset of the parameter space.

Multiplication by the fixed product of the \(m=d-3\) hyperplane equations gives a morphism \(S\to\mathcal P_d\). The image lies in \(W\) by \eqref{eq:cubic-5-5} and the Frobenius invertibility just proved. On \(W\), Section~\ref{sec:pn-2} constructs the finite étale rank-\(r\) local system \(\mathcal L=\mathcal M^{F=1}\). Its restriction to \(S\) is exactly the boundary local system of \eqref{eq:cubic-6-1}.

Let \(W^{\mathrm{sm}}\subset W\) be the open of smooth hypersurfaces. It is nonempty by Section~\ref{sec:pn-3}. On \(W^{\mathrm{sm}}\), the Artin--Schreier sequence

\[
0\to\mathbf F_p\to\mathcal O\xrightarrow{F-1}\mathcal O\to0
\]

identifies \(\mathcal L\) with the top étale cohomology local system. Indeed the lower positive coherent cohomology groups vanish. For \(n=2\), the preceding term is \(H^0(\mathcal O)=\kappa\), on which \(F-1\) is surjective; relatively this is the usual surjective Artin--Schreier map on the étale site. The coherent/étale comparison and the explicit higher-direct-image calculation give the same identification in families. The fibres therefore are \(H^{n-1}_{\mathrm{et}}(Y_{\bar s},\mathbf F_p)\), with their natural monodromy.

The homomorphism \(\pi_1(W^{\mathrm{sm}})\to\pi_1(W)\) is surjective. To check it, every connected finite étale cover of the normal integral scheme \(W\) is itself integral; its restriction to a nonempty open is still connected. This is the connected-cover criterion for surjectivity, with the normal-cover assertion given explicitly in \citet[Lemma~58.11.2, Tag~0BQL]{StacksNormal}. Thus the monodromy image on \(\mathcal L\) over \(W^{\mathrm{sm}}\) is the same as over \(W\), and contains the boundary image coming from \(S\).

The ranks agree by \eqref{eq:cubic-5-5}. The comparison is realized by an actual finite étale local system on \(W\), without extending the generic system across a rank-drop locus.

\subsection{All traces in the geometric image}\label{sec:cubic-traces}

Because \(m\ge n-1\), there is at least one node factor and at least one elliptic factor. Theorem~\ref{thm:joint} lets us vary those two selected factors independently while fixing all other listed factors. Any additional extension data in the boundary representation do not affect trace, by \eqref{eq:cubic-6-1}.

On \(\mathbf F_p^{Z_J}/\mathbf F_p\mathbf1\), the three conjugacy types have traces

\begin{equation}
\begin{array}{c|ccc}
\text{type}&1&\text{transposition}&\text{3-cycle}\\\hline
\text{trace}&2&0&-1.
\end{array}
\label{eq:cubic-8-1}
\end{equation}

For odd \(p\), let the selected elliptic scalar range through \(\mathbf F_p^\times\). Combining it with the identity and transposition in the selected node factor gives, up to a fixed common translation, the union

\[
\mathbf F_p^\times\ \cup\ (2+\mathbf F_p^\times)
=\mathbf F_p.
\]

The equality holds because the two sets omit respectively 0 and 2, which are distinct in odd characteristic. For \(p=2\), the elliptic scalars are all 1, and the identity and 3-cycle in the selected node factor have traces 0 and 1. They already give every element of \(\mathbf F_2\), again up to a fixed translation. Consequently the boundary traces, and hence the geometric generic traces by Section~\ref{sec:cubic-generic}, exhaust \(\mathbf F_p\).
\subsection{Determinants and semisimple rank-one elements}
\begin{proof}[Proof of Proposition~\ref{prop:determinant-reflections}]
For $d=n+1$, the determinant assertion is the rank-one image calculation in Section~\ref{sec:pn-5}. For $p=2$, it is automatic because $\mathbf F_2^\times=\{1\}$. We may therefore assume $p$ odd and $d\ge n+2$, and put $m=d-3\ge n-1$.

The boundary representation of Sections~\ref{sec:cubic-cech}--\ref{sec:cubic-generic} has dimension $h$ and image contained, after identifying fibres, in the geometric generic image. Its invariant filtration has constant pieces, the node quotients $\mathbf F_p^{Z_J}/\mathbf F_p\mathbf1$, and the elliptic lines. Determinants multiply across an invariant filtration. The determinant on a node quotient is the sign of the permutation: the full permutation module has that determinant, and its diagonal constant line has determinant $1$. Consequently a boundary element with node permutations $\sigma_J$ and elliptic scalars $a_I$ has determinant
\[
\prod_J\operatorname{sgn}(\sigma_J)\prod_I a_I.
\]
This uses no splitting, including in characteristic $3$.

There is at least one elliptic factor. By Theorem~\ref{thm:joint}, for each $a\in\mathbf F_p^\times$ there is a boundary element $g$ acting by $a$ on one selected elliptic line, by $1$ on every other elliptic line, and trivially on all node factors. It acts trivially on the constant pieces. Thus $\det(g)=a$, proving geometric determinant surjectivity. The arithmetic image contains the geometric one and has determinant in the same group $\mathbf F_p^\times$, so its determinant image is also full.

Take now $a\ne1$. The characteristic polynomial of this same filtered automorphism is
\[
(T-a)(T-1)^{h-1}.
\]
Its primary decomposition is defined over $\mathbf F_p$: there is a one-dimensional $a$-eigenspace and a generalized $1$-eigenspace of dimension $h-1$. On the latter write $g=I+N$, with $N^h=0$. Choose $r$ with $p^r\ge h$. In characteristic $p$,
\[
(I+N)^{p^r}=I+N^{p^r}=I,
\qquad a^{p^r}=a.
\]
Hence $g^{p^r}$ is diagonalizable over $\mathbf F_p$, with eigenvalue $a$ once and eigenvalue $1$ with multiplicity $h-1$. It belongs to the same geometric group and is conjugate to the displayed diagonal matrix. This proves the claimed rank-one semisimple elements, without an irreducibility or full-image assertion.
\end{proof}

\section{Every residue in every parameter open}\label{sec:pn-9}
We now prove the finite-cover point-existence statement and apply it to the geometric trace calculation. The only point-existence theorem used is the following consequence of Lang--Weil: a nonempty geometrically integral variety over $\mathbf F_q$ has a point over $\mathbf F_{q^r}$ for every sufficiently large $r$ \citep{LangWeil}; the precise extension-field formulation is \citet[Theorem~7.7.1(iii), pp.~222--223]{Poonen}.

\begin{lemma}[Realizing a finite geometric image]\label{lem:finite-image-twist}
Let $U$ be a smooth geometrically integral variety over a finite field $k$, and let $\mathcal V$ be a finite-rank $\mathbf F_p$ local system on $U$. Write $G$ for its finite arithmetic image and $H$ for its geometric image. After a finite extension of $k$, every conjugacy class of $H$ is realized as an arithmetic Frobenius class over a further finite extension, in every nonempty open of $U$.
\end{lemma}
\begin{proof}
The geometric image is normal in $G$, and the finite quotient $G/H$ is a quotient of the absolute Galois group of the finite field. Replace $k$ by the finite extension that kills this quotient. The arithmetic image is then $H$, while the geometric image remains $H$. The connected $H$-torsor $Z\to U$ associated with the representation is geometrically connected. Since it is finite \'etale over the smooth base, it is geometrically integral. A further restriction to any prescribed nonempty open preserves this property. Shrink that open to an affine smooth open.

Write $k=\mathbf F_q$ and fix $g\in H$ of order $s$. Twist the descent action of arithmetic Frobenius on $Z_{\overline{\mathbf F_q}}$ by right multiplication by $g^{-1}$. The torsor group is constant, so these two actions commute. This descent datum is effective: over $\mathbf F_{q^s}$ its $s$-fold action is the original one, and affine finite Galois descent gives a variety $Z_g$ over $\mathbf F_q$. Its geometric base change is isomorphic to that of $Z$, hence is integral.

A point of $Z_g(\mathbf F_{q^r})$ corresponds to a point $z$ of $Z(\overline{\mathbf F_q})$ satisfying
\begin{equation}
\operatorname{Frob}_{q,\mathrm{arith}}^r(z)=zg^r.
\label{eq:finite-twist-frobenius}
\end{equation}
Lang--Weil supplies such a point for every sufficiently large $r$. Choose one with $r\equiv1\pmod s$. Its image in $U(\mathbf F_{q^r})$ has arithmetic Frobenius conjugate to $g$. This is precisely the required realization after finite extensions.
\end{proof}

\begin{proof}[Proof of Theorem~\ref{thm:threshold}]
Nonemptiness is Section~\ref{sec:pn-3}, and the sharp lower obstruction is Section~\ref{sec:pn-4}. Suppose $d\ge n+2$ and put $m=d-3$. Over an algebraic closure, choose $m$ hyperplanes in general position. They and the associated family descend to some finite extension of the original finite field. Sections~\ref{sec:cubic-joint}--\ref{sec:cubic-generic} show that the geometric generic image has every trace in $\mathbf F_p$.

Let $U\subset S_d$ be the prescribed nonempty open. Its intersection with the smooth ordinary open is nonempty by geometric integrality. Restriction to this smaller open leaves the geometric image unchanged: a connected finite \'etale cover of a normal integral scheme remains connected on a nonempty open, as used in Section~\ref{sec:cubic-generic}. For a prescribed residue $a$, choose an element $h$ in the geometric image with
\[
\operatorname{tr}(h)=(-1)^{n-1}(a-1).
\]
Apply Lemma~\ref{lem:finite-image-twist} to the arithmetic conjugacy class of $h^{-1}$. At the resulting point, geometric Frobenius has the class of $h$. The convention established in \eqref{eq:pn-2-4} and the point-count formula \eqref{eq:pn-4-2} give
\[
\#Y(k')\equiv1+(-1)^{n-1}\operatorname{tr}(h)=a\pmod p.
\]
Composing the finite extensions returns exactly a finite extension of the originally given field. The constructed hypersurface lies in the original $U$ and is smooth. This proves the theorem in every characteristic and every degree in the stated range.
\end{proof}

\section{The characteristic-two covector, with global descent}\label{sec:pn-6}

This section constructs a global invariant covector and proves that it never vanishes on the ordinary open.

Suppose now that \(p=2\) and

\begin{equation}
                         d=n+1+2e,\qquad e\geq0.
\end{equation}

Let \(S_j\) denote homogeneous polynomials of degree \(j\) in the
\(n+1\) variables. The dual of the coherent space in \eqref{eq:pn-1-3} is
\(S_{2e}\), with \(x^{-a}\) paired with \(x^{a-\mathbf1}\).
Introduce the differential operator and homogeneous polynomial

\begin{equation}
 D=\partial_{x_0}\partial_{x_1}\cdots\partial_{x_n},\qquad
                         \xi_f=Df\in S_{2e}.                           \label{eq:pn-6-1}
\end{equation}

Every exponent occurring in \(\xi_f\) is even: a monomial of \(f\)
survives \(D\) exactly when each of its exponents is odd. In particular,
every first derivative of \(\xi_f\) is zero.

\subsection{The dual semilinear formula}

Over a perfect field, the Frobenius-dual Cartier map on \(S_{2e}\) is

\begin{equation}
                         C_f(g)=\bigl(D(fg)\bigr)^{1/2}.               \label{eq:pn-6-2}
\end{equation}

Here the square root means halving even monomial exponents and taking
the unique square roots of their coefficients. To verify the formula
from \eqref{eq:pn-1-4}, write \(g=\sum_b g_bx^{b-\mathbf1}\). The coefficient of
\(x^{2(a-\mathbf1)}\) in \(D(fg)\) is

\begin{equation}
                  \sum_{b\in I_d}f_{2a-b}g_b=(B_2^tg)_a.              \label{eq:pn-6-3}
\end{equation}

Thus \eqref{eq:pn-6-2} is precisely \(g\mapsto(B_2^tg)^{[1/2]}\), and its defining
duality is

\begin{equation}
                     \langle Fv,g\rangle=\langle v,C_f(g)\rangle^2.    \label{eq:pn-6-4}
\end{equation}

Since every derivative of \(\xi_f\) vanishes, the product rule gives

\begin{equation}
                 D(f\xi_f)=\xi_fDf=\xi_f^2,
 \qquad                         C_f(\xi_f)=\xi_f.                     \label{eq:pn-6-5}
\end{equation}

There is no perfection assumption on the universal identity underlying
\eqref{eq:pn-6-5}. If \(\xi\) is its coefficient column in \(S_{2e}\), \eqref{eq:pn-6-3} gives
the polynomial identity

\begin{equation}
                             B_2^t\xi=\xi^{[2]}.                      \label{eq:pn-6-6}
\end{equation}

For an alternative direct check of its coefficients, \(\xi_b=f_b\)
when every coordinate of \(b\) is odd, and \(\xi_b=0\) otherwise.
The summands of \(\sum_b f_{2a-b}\xi_b\) cancel in pairs under
\(b\leftrightarrow2a-b\). A surviving exponent \(2a-b\) is again
positive and odd; negative exponents contribute zero. The only possible
fixed term is \(b=a\), yielding exactly \(\xi_a^2\).

\subsection{Nonvanishing on the entire ordinary open}

Assume \(B_2\) is invertible over an algebraically closed field. If
\(\xi_f=0\), take any nonzero \(h\in S_e\). All derivatives of \(h^2\)
vanish, so

\begin{equation}
                 C_f(h^2)=\bigl(h^2Df\bigr)^{1/2}=0.
\end{equation}

This contradicts invertibility of \(C_f\), because \(h^2\neq0\).
Thus \(\xi_f\neq0\) at every geometric ordinary point. This proves
nowhere-vanishing, not merely generic nonvanishing.

If \(v\) is a Lang fixed vector, \eqref{eq:pn-6-6} gives

\begin{equation}
 (\xi^tv)^2=(\xi^{[2]})^tv^{[2]}
           =\xi^t B_2v^{[2]}=\xi^tv.
\end{equation}

Consequently \(v\mapsto\xi^tv\) takes values in \(\mathbf F_2\).
It is nonzero on the fixed-vector space because that space spans the
coherent vector space by Section~\ref{sec:pn-2} and \(\xi\neq0\).

\subsection{Descent, coordinate independence, and the invariant hyperplane}

The identity \eqref{eq:pn-6-6} has coefficients in \(\mathbf F_2\); it therefore defines
a morphism of local systems to the constant sheaf \(\mathbf F_2\), rather
than just a fiberwise collection of functionals. It respects changes of
equation. Namely, if \(f'=uf\), then \(\xi_{f'}=u\xi_f\) while coherent
coordinates change by \(v'=u^{-1}v\). The pairing is unchanged.

Equivalently, under adjunction this dual section is the residue of

\begin{equation}
                 \frac{\xi_f\Omega}{f},\qquad
 \Omega=\sum_i(-1)^i x_i\,dx_0\wedge\cdots\widehat{dx_i}\cdots\wedge dx_n.
\end{equation}

The expression is homogeneous of degree zero and is invariant under
scaling \(f\). It is also compatible with linear coordinate changes.
If \(y=Ax\), the chain rule and \(\partial_{y_j}^2=0\) in characteristic
two imply

\begin{equation}
 D_x(f(Ax))=\det(A)(D_yf)(Ax).
\end{equation}

Indeed only permutations of the distinct derivatives survive; their
permanent is the determinant in characteristic two. Also
\(A^*\Omega_y=\det(A)\Omega_x\), giving the claimed compatibility.

Thus the construction supplies exactly a nowhere-zero invariant
\(\mathbf F_2\)-valued covector \(\lambda\), defined over \(\mathbf F_2\),
on the universal ordinary local system. Both geometric and arithmetic
images lie in \(\operatorname{Stab}(\lambda)\). If \(d\geq n+3\), then
\(h>1\), so \(\ker\lambda\) is a proper nonzero invariant hyperplane.
This proves the canonical-covector assertion and its reducibility consequence. By itself this covector argument does not prove equality with the whole stabilizer. Equality is established by the plane and higher-dimensional group arguments below.

\section{Full monodromy for \(p=2,d=n+2\)}\label{sec:pn-7}

This section determines the full monodromy group in the first characteristic-two degree above the obstruction.

In this degree,

\begin{equation}
 I_d=\{\mathbf1+e_j:0\leq j\leq n\},\qquad h=n+1.
\end{equation}

Equation \eqref{eq:pn-1-4} becomes

\begin{equation}
 (B_2)_{i,j}=f_{\mathbf1+2e_j-e_i}.                                    \label{eq:pn-7-1}
\end{equation}

These \(h^2\) coefficients are all distinct. The diagonal exponents have
one entry 2 and all other entries 1. An off-diagonal exponent has entry
0 at \(i\), entry 3 at \(j\), and all other entries 1. None is the pure
power exponent \((d,0,\ldots,0)\), since \(n\geq2\) and \(d\geq4\).
On the projective coefficient chart in which the coefficient of \(x_0^d\)
is 1, the entries of \(B_2\) are therefore independent affine coordinates.
Write the remaining coefficient coordinates as \(z_1,\ldots,z_s\).

On the ordinary part of this chart, the parameter space before imposing
smoothness is

\begin{equation}
              X=\mathrm{GL}_h\times\mathbf A^s,
\end{equation}

with the first factor parametrizing the matrix \(B_2\).
The basis torsor is the incidence

\begin{equation}
       \{(B,T,z):T=BT^{[2]},\ B,T\in\mathrm{GL}_h\}.
\end{equation}

Solving for \(B\) identifies its total space with

\begin{equation}
 \mathrm{GL}_h\times\mathbf A^s,
 \qquad (T,z)\longmapsto\bigl(T(T^{[2]})^{-1},z\bigr).                 \label{eq:pn-7-2}
\end{equation}

This is the finite étale \(\mathrm{GL}_h(\mathbf F_2)\)-torsor of Section~\ref{sec:pn-2}.
Its total space is geometrically integral. The smooth locus is a nonempty
open by Section~\ref{sec:pn-3}, and its inverse image under \eqref{eq:pn-7-2} is a nonempty open
of a geometrically integral scheme, hence remains connected. Therefore
the geometric monodromy on the smooth ordinary locus is
\(\mathrm{GL}_h(\mathbf F_2)\). The arithmetic image contains it and lies
in the same group, so it is equal to it. This works over every perfect
field of characteristic two.

This proof determines the full image in the first degree above the obstruction. It does not make the higher-degree Hasse--Witt matrix a matrix of independent
coefficients.

\section{A direct matrix realization when \(p=2,d=n+2\)}\label{sec:pn-8}

Here the rational description \eqref{eq:pn-7-2} permits a direct proof of the point
existence needed for residues. This section does not use Lang--Weil or a
finite-field Chebotarev theorem.

Let \(k=\mathbf F_q\), with \(q\) a power of two, and let \(U\) be any
nonempty open in the smooth degree-\(d\) parameter space. Its intersection
with the chart and ordinary locus of Section~\ref{sec:pn-7} is nonempty. Shrinking
within that intersection, it contains a nonempty principal open

\begin{equation}
       F(B,z)\neq0\quad\text{inside }\mathrm{GL}_h\times\mathbf A^s,     \label{eq:pn-8-1}
\end{equation}

where \(F\) is a polynomial over \(\mathbf F_q\). Denominators that are
powers of \(\det B\) may be cleared.

Fix any \(g\in\mathrm{GL}_h(\mathbf F_2)\), and choose a finite extension
\(\mathbf F_Q\supseteq\mathbf F_q\). By the prime-power fixed-basis
argument in Section~\ref{sec:pn-2}, there is \(C\in\mathrm{GL}_h(\overline{\mathbf F_q})\)
such that

\begin{equation}
                         C^{[Q]}=Cg.                                  \label{eq:pn-8-2}
\end{equation}

For example, apply that argument to \(g^{-t}\) and transpose the resulting
fixed-basis matrix. Put \(B_0=C(C^{[2]})^{-1}\). Since \(g^{[2]}=g\),
\eqref{eq:pn-8-2} implies \(B_0^{[Q]}=B_0\), so \(B_0\in\mathrm{GL}_h(\mathbf F_Q)\).

For \(S\in\mathrm{GL}_h(\mathbf F_Q)\), define

\begin{equation}
 T=SC,\qquad B(S)=SB_0(S^{[2]})^{-1}.                                  \label{eq:pn-8-3}
\end{equation}

Then \(B(S)\) is defined over \(\mathbf F_Q\),
\(B(S)T^{[2]}=T\), and \(T^{[Q]}=Tg\). Thus arithmetic Frobenius on the
Lang basis has exactly matrix \(g\).

Over the algebraic closure, the map \(S\mapsto B(S)\) is the finite étale
surjective map \eqref{eq:pn-7-2} after the invertible substitution \(T=SC\).
Consequently \(F(B(S),z)\) is a nonzero rational function. Multiplying
by a sufficient power of \(\det S\) clears its denominators and gives a
nonzero polynomial \(G_Q(S,z)\) over \(\mathbf F_Q\). The total degree of
\(G_Q\) has a bound depending only on \(F\) and \(h\), not on \(Q\) or the
entries of \(B_0\). This follows directly from the cofactor formula for
\((S^{[2]})^{-1}\). The nonzero polynomial
\((\det S)G_Q(S,z)\) has the same kind of uniform degree bound.

A nonzero polynomial of total degree \(D\) in \(r\) variables over
\(\mathbf F_Q\) vanishes at at most \(D Q^{r-1}\) points. This elementary
bound follows by induction on \(r\), writing the polynomial in the last
variable and separating points where its leading coefficient vanishes.
Thus, for every sufficiently large \(Q\), there are \(S,z\) over
\(\mathbf F_Q\) with \(\det S\neq0\) and \(F(B(S),z)\neq0\).
They give a point of the prescribed open \(U\) with arithmetic Frobenius
matrix \(g\).

The traces of elements of \(\mathrm{GL}_h(\mathbf F_2)\), with \(h=n+1\geq3\),
take both values in \(\mathbf F_2\). Explicitly, the companion matrix of
\(X^h-tX^{h-1}+1\) is invertible and has trace \(t\), for either \(t\).
Choose \(g\) as its inverse. Equations \eqref{eq:pn-2-4} and \eqref{eq:pn-4-2} then give residue
\(1+t\). Both residues occur over finite extensions in every such \(U\).

\section{Lifted divisors of an abelian curve cover}\label{sec:plane-wreath}
This section determines the node permutations together with their lifts to a finite abelian cover of the curve.

\begin{lemma}\label{lem:plane-wreath}
Let $C$ be a smooth projective connected curve over an algebraically closed field, and let $L$ be very ample of degree $N\ge2$. Suppose a general nonreduced member of $|L|$ has one point of length two and $N-2$ distinct simple points. Let $C'\to C$ be a connected finite \'etale torsor for a finite abelian group $G$, written additively. On the open $B^{\mathrm{red}}\subset B=\mathbf P H^0(C,L)$ of reduced divisors, the monodromy on all their lifted points is, after choices of labels,
\[
 G_0^N\rtimes S_N,\qquad
 G_0^N=\{(g_1,\ldots,g_N)\in G^N:\ \textstyle\sum_i g_i=0\}.
\]
Here the group acts on $G\times\{1,\ldots,N\}$ by translations within the blocks and permutations of the blocks.
\end{lemma}

\begin{proof}
The incidence space of a lifted point $Q\in C'$ and a section vanishing at its image is a projective bundle over $C'$, with fibre $\mathbf P H^0(C,L(-\pi(Q)))$. It is irreducible. Its open above $B^{\mathrm{red}}$ is a connected finite \'etale cover. Thus its monodromy $H$ is transitive on $G\times\{1,\ldots,N\}$. Deck translations commute with monodromy, so $H\subset G^N\rtimes S_N$.

The tangent-divisor incidence over $C$ is a projective bundle because the embedding separates first jets. Its image is the irreducible nonreduced-divisor discriminant. It is a hypersurface: each divisor has finitely many points, so the tangent incidence has finite generic fibre. At its generic geometric point, the hypothesis gives one double point and all other points simple.

After strict henselization at this discriminant, the point-incidence cover has one degree-two branch and $N-2$ \'etale branches. The \'etale cover $C'\to C$ splits at the colliding point and at the other points. Its inertia permutation therefore has the form
\[
 t_{ij}(g):(x,i)\longmapsto(x+g,j),\qquad
 (x,j)\longmapsto(x-g,i),
\]
and fixes all other blocks. This remains valid in characteristic two. The local generic extension has separable degree two and Galois group of order two. A purely inseparable residue extension at the generic tangent point does not add a translation: after separable residue closure, the \'etale cover above that point is split, and only the degree-two separable branch extension remains.

Normalize $B$ in the Galois closure of the point-incidence extension. The quotient by the normal subgroup generated by all divisorial inertia is unramified in codimension one. It is a finite normalization in a separable function field, so its source is normal, every component dominates the regular target $B$, and the required local dimensions agree. Purity of branch locus makes this quotient \'etale \citep[Lemma~58.21.4]{StacksPurity}. Projective space over an algebraically closed field has no connected nontrivial finite \'etale cover. Hence $H$ is generated by the displayed twisted transpositions and their conjugates.

Here one may verify the last projective-space assertion directly. Riemann--Hurwitz excludes a connected nontrivial finite \'etale cover of $\mathbf P^1$ \citep[Section~53.12, Lemma~53.12.2]{StacksRH}. For $\mathbf P^r$, $r\ge2$, pull a cover back to the blowup of a point, which is a $\mathbf P^1$-bundle over $\mathbf P^{r-1}$. It splits on each fibre. Proper cohomology and base change \citep[Lemma~30.22.1]{StacksBaseChange} for its finite locally free algebra descend it to a finite \'etale algebra on $\mathbf P^{r-1}$: fibrewise $H^0$ has the constant degree, $H^1$ vanishes, and the evaluation map is an isomorphism. The exceptional section maps to the original one point; restricting there shows the descended algebra is trivial. The pulled-back cover is therefore trivial. A connected cover remains connected after this birational base change, so its degree is one.

It remains to identify the generated subgroup. Form the graph of the transpositions in a generating set. Its projection is transitive, so the graph is connected. Choose a spanning tree. Changing the origins in the $G$-labelled blocks makes each tree generator $t_{ij}(0)$. These generate the ordinary $S_N$. Multiplying any further $t_{ij}(g)$ by $t_{ij}(0)$ gives a translation with entries $g,-g$ in two coordinates. Let $G'\subset G$ be generated by the labels that occur. Conjugating by $S_N$ gives every pair translation with entries in $G'$, hence the whole zero-sum subgroup $(G')_0^N$. Conversely all generators lie in $(G')_0^N\rtimes S_N$, so this is exactly $H$. Its orbits on $G\times\{1,\ldots,N\}$ are the cosets of $G'$ in $G$. Transitivity forces $G'=G$.
\end{proof}

The zero-sum condition is essential. This is the monodromy of one abelian curve cover lifted over all points of a divisor, rather than unrelated covers of the individual points.

\section{Simple contacts for the line and conic series}\label{sec:plane-contact}
This section checks the geometric hypothesis of Lemma~\ref{lem:plane-wreath} for the plane curves used below.

Let $C\subset\mathbf P^2$ be a general smooth curve of degree $e\ge3$, and let $q=1$ or $2$. Since $q<e$ and $H^1(\mathbf P^2,\mathcal O(q-e))=0$, restriction identifies $H^0(C,\mathcal O_C(q))$ with all degree-$q$ plane forms. Thus its complete series is precisely the moving-line or moving-conic series, and it is very ample.

Consider triples $(C,R,P)$ with $R$ a smooth line or conic, $P\in R$, and $C|_R$ vanishing to order at least two at $P$. For fixed $(R,P)$ these are two independent linear conditions on the equation of $C$. The total incidence is irreducible: it is a projective bundle over the irreducible family of pairs $(R,P)$, with suitable nonempty opens removed. It dominates the parameter space of $C$, since at every smooth point tangent lines, and tangent smooth conics, exist.

The incidence has a nonempty open where $C$ is smooth, the contact at $P$ has length exactly two, and all other intersections are transverse. Fix $R,P$ to check this. Restrictions of degree-$e$ plane equations give the complete binary series of degree $qe\ge3$. It separates a length-two condition at $P$ together with a length-two condition at $Q\ne P$. A second double point imposes two conditions while $Q$ varies in a curve, giving a proper closed locus. Exact order two at $P$ is also a nonempty open. For smoothness away from $P$, choose a linear form $\ell$ vanishing at $P$ but not at a given $Q$. Products $\ell^2g$, with $\deg g=e-2\ge1$, satisfy the imposed double condition and separate first jets at $Q$. The usual singular-point incidence count applies. At $P$, prescribe a nonzero differential normal to $R$. These requirements are compatible, including in characteristic two.

Consequently the good open dominates a dense open of the $C$-parameter space. For such a fixed $C$, its tangent incidence is irreducible and contains a simple collision, so its generic tangent divisor has exactly the form required in Lemma~\ref{lem:plane-wreath}. Generic Hasse--Witt invertibility is another nonempty open, by the generic-invertibility theorem recalled earlier. We choose $C$ in the intersection of these opens. Finally, restricting the complete line or conic series to the nonempty open of smooth $R$ preserves its finite monodromy image: a connected finite \'etale cover of the normal integral parameter space stays connected after a nonempty open restriction, as in Section~\ref{sec:cubic-generic}.

\section{Actual transvections from nodal gluing}\label{sec:plane-transvections}
This section obtains unipotent elements on the entire boundary cohomology, beyond their actions on its graded pieces.

Fix an ordinary $C$ of degree $e\ge3$ as above, and let $R$ vary among smooth rational curves of degree $q\in\{1,2\}$ transverse to $C$. Put
\[
 N=eq,\quad Y=C\cup R,\quad U=H^1_{\mathrm{et}}(C,\mathbf F_p),
 \quad P=\mathbf F_p^N/\mathbf F_p\mathbf1.
\]
Normalization gives an equivariant exact sequence
\begin{equation}\label{eq:plane-extension}
 0\longrightarrow P\longrightarrow H^1_{\mathrm{et}}(Y,\mathbf F_p)
                   \longrightarrow U\longrightarrow0.
\end{equation}
The same coherent sequence has bijective Frobenius on its graph part and on $H^1(C,\mathcal O_C)$. Thus all coherent $H^1(Y,\mathcal O_Y)$ is stable, of rank $\dim U+N-1$, the arithmetic genus of the degree-$d=e+q$ plane divisor. The finite boundary Lang part consequently has the full generic rank. The boundary specialization argument of Section~\ref{sec:cubic-generic} places this representation inside generic ordinary geometric monodromy.

Take the connected maximal elementary abelian \'etale cover $C'\to C$ with group $G=U^\vee$. Choosing lifts of the $N$ nodes identifies \eqref{eq:plane-extension} with $P\oplus U$ as vector spaces. Changing these choices by $(g_1,\ldots,g_N)\in G^N$ changes the splitting by
\begin{equation}\label{eq:plane-gluing-translation}
 \alpha\longmapsto[(\alpha(g_1),\ldots,\alpha(g_N))]\in P.
\end{equation}
To verify this on actual torsors, glue an $\mathbf F_p$-torsor on $C$ representing $\alpha$ to the trivial torsor on $R$, with fibre identifications at all nodes. Changing the trivialization on $R$, or translating the torsor on $C$, adds a common constant, which is exactly the diagonal factored out in $P$. Pushouts of the single $G$-torsor make these choices simultaneously and linearly in $\alpha$. The glued cover is finite \'etale: at every node, matched pairs of branches have completed local rings that are products of copies of $k[[x,y]]/(xy)$; away from nodes the covers are \'etale on the two components.

Lemma~\ref{lem:plane-wreath} gives the lifted-node group $G_0^N\rtimes S_N$. If $p\nmid N$, its translations map through \eqref{eq:plane-gluing-translation} onto all of $\operatorname{Hom}(U,P)$, since averaging $N^{-1}\sum g_i$ identifies $G_0^N$ with $G^N/\operatorname{diag}G$. If $p\mid N$, their image is exactly
\[
 \operatorname{Hom}(U,P_0),\qquad
 P_0=\ker\bigl(\textstyle\sum:P\longrightarrow\mathbf F_p\bigr).
\]
The sum is well-defined in the latter case because $N=0$ in $\mathbf F_p$. Thus boundary monodromy contains the actual groups
\begin{equation}\label{eq:plane-radicals}
 \begin{cases}
 I+\operatorname{Hom}(U,P),&p\nmid N,\\
 I+\operatorname{Hom}(U,P_0),&p\mid N.
 \end{cases}
\end{equation}
These endomorphisms factor through the quotient $U$ and the subspace $P$, so they square to zero. The conclusion concerns the extension itself, not only its associated graded representation.

\section{A common marking at a line arrangement}\label{sec:plane-marking}
This section identifies all boundary subspaces in one fibre so that their transvections can be combined.

Let $Y_0$ be a union of $d$ general labelled lines, with no triple point. Its cohomology is
\[
 V=\mathbf F_p^{E(K_d)}/\delta\mathbf F_p^{V(K_d)},\qquad
 \dim V=g=\frac{(d-1)(d-2)}2.
\]
Write $e_{ij}=-e_{ji}$ for oriented edge classes. The star relations are
$\sum_{j\ne i}e_{ij}=0$. The ordered line-configuration space is connected, and its unordered quotient gives the full $S_d$ action on this graph cohomology. The arrangement is ordinary because its entire $H^1$ is the stable graph part. It is therefore a valid fibre of the Lang system on the full ordinary equation open.

Specialize a degree-$(d-1)$ curve $C$ to all lines except line $i$, and take $R$ to be line $i$. In the connected open family of pairs $(C,R)$, the curve $C$ can be smooth or this ordinary nodal arrangement; intersections are \'etale and avoid internal nodes. Coherent $H^1$ remains a vector bundle with invertible Frobenius. Consequently \eqref{eq:plane-extension} is a sequence of finite local systems across the specialization. Transport to $Y_0$ identifies its kernel with
\[
 P_i=\operatorname{span}\{e_{ij}:j\ne i\},\qquad \dim P_i=d-2.
\]
When $p\nmid d-1$, the transported group in \eqref{eq:plane-radicals} is the whole $I+\operatorname{Hom}(V/P_i,P_i)$. A different marking preserving $P_i$ preserves this entire subgroup. If $p\mid d-1$, its image is $I+\operatorname{Hom}(V/P_i,P_{i,0})$, where $P_{i,0}$ is the node sum-zero subspace.

For a moving conic, specialize $C$ to $d-2$ lines and the conic to the remaining two lines. This reducible conic still has $H^1=0$. The kernel now identifies with the cut subspace
\[
 P_B=\operatorname{span}\{e_{ij}:i\in B,\ j\notin B\},\qquad
 |B|=2,\quad \dim P_B=2d-5.
\]
It is the kernel of restriction to the complementary complete graph on $d-2$ vertices, since $H^1(K_2)=0$. Conjugating by $S_d$ supplies the corresponding subgroups for every vertex and every two-vertex cut.

\section{Generating linear groups from graph subspaces}\label{sec:plane-graph}
This section proves the linear algebra needed to pass from the boundary transvections to special linear groups.

\begin{lemma}\label{lem:plane-stars}
For $d\ge4$ and $i\ne j$, one has $P_i\cap P_j=\mathbf F_p e_{ij}$. A subgroup of $\mathrm{GL}(V)$ containing $I+\operatorname{Hom}(V/P_i,P_i)$ for every vertex contains $\mathrm{SL}(V)$.
\end{lemma}
\begin{proof}
Suppose cochains supported on stars $i$ and $j$ differ by $\delta f$. On all edges among the remaining $d-2$ vertices that difference is zero, so $f$ is constant there. Subtract the constant. The values $f_i,f_j$ then permit subtraction of star coboundaries, leaving only edge $ij$. Its class is nonzero because it lies on a triangle. This proves the intersection assertion.

For $v=e_{ij}$, the group contains $I+v\phi$ for every $\phi\in\operatorname{Ann}(P_i)$ or $\operatorname{Ann}(P_j)$. These transvections commute and add their functionals, and
\[
 \operatorname{Ann}(P_i)+\operatorname{Ann}(P_j)=\operatorname{Ann}(v).
\]
Thus every transvection of centre $v$ occurs. Edge classes span $V$, so a basis can be selected from them. The group contains all elementary matrices in that basis, which generate $\mathrm{SL}(V)$.
\end{proof}

\begin{lemma}\label{lem:plane-cuts}
For $d\ge6$, a subgroup containing $I+\operatorname{Hom}(V/P_B,P_B)$ for every two-vertex set $B$ contains $\mathrm{SL}(V)$.
\end{lemma}
\begin{proof}
Fix edge $ij$. The intersection of $P_B$ over cuts separating $i$ and $j$ is $\mathbf F_p e_{ij}$. Indeed a class in that intersection is zero on the cycle space of every complementary $K_{d-2}$. Every triangle not using edge $ij$ lies in such a complement. If it contains $i$ but not $j$, omit $j$ and another vertex outside it. If it contains neither, omit $i$ and a vertex outside the triangle and $i,j$, using $d\ge6$. The other cases are symmetric.

The cycle space of $K_d$ with edge $ij$ removed is generated by those triangles. Cycles of length at least four have a chord; if one possible chord is the removed edge, the other chord of that four-cycle is present. Hence the class vanishes on every cycle of the connected edge-deleted graph. Subtracting a coboundary makes it supported only on $ij$, proving the intersection assertion.

The sum of the annihilators of these cuts is therefore $\operatorname{Ann}(e_{ij})$. Products of their common-centre transvections supply every functional in this annihilator. Selecting an edge basis gives all elementary matrices, as in Lemma~\ref{lem:plane-stars}.
\end{proof}

\section{All plane degrees in odd characteristic}\label{sec:plane-odd}
This section proves the special linear lower bound in odd characteristic, including the exceptional quartic over characteristic three.

If $p\nmid d-1$, use a fixed ordinary curve of degree $d-1$ and a moving line. Equations~\eqref{eq:plane-radicals} and Lemma~\ref{lem:plane-stars} give $\mathrm{SL}_g(\mathbf F_p)$.

If $p\mid d-1$, use a fixed ordinary curve of degree $d-2$ and a moving conic. The node number $2(d-2)$ is prime to the odd prime $p$. For $d\ge6$, Lemma~\ref{lem:plane-cuts} gives $\mathrm{SL}_g(\mathbf F_p)$. The only omitted case with $d\ge4$ is $d=4,p=3$; degree five has $d-1=4$, with no odd prime divisor.

For this exceptional case, use the basis $x=e_{12},y=e_{13},z=e_{23}$ of $H^1(K_4,\mathbf F_3)$. The line--cubic boundary has three nodes. Its sum-zero translations give all scalar multiples of
\[
 \begin{array}{c|c|c}
 &w_i&\phi_i\\\hline
 T_1=I+w_1\phi_1&(1,-1,0)&z^*\\
 T_2=I+w_2\phi_2&(-1,0,-1)&y^*\\
 T_3=I+w_3\phi_3&(0,-1,1)&x^*.
 \end{array}
\]
Each $w_i$ is the difference of two outgoing edge classes, and $\phi_i$ annihilates its star. Their scalings are arbitrary because the elliptic quotient has dimension one. The three $w_i$ form a basis: their determinant is $-2=1$ in $\mathbf F_3$.

The elements $T_2,T_3$ generate $\mathrm{SL}$ on $\operatorname{span}(w_2,w_3)$ and fix the complementary line $\mathbf F_3z$. In the indicated plane basis they are the two elementary matrices with nonzero coefficient $-1$, whose powers supply all coefficients. Put
\[
 u=-w_2+w_3=(1,-1,-1),\qquad \psi=x^*+y^*.
\]
The functional $\psi$ kills $u$ and $z$, so $S=I+u\psi$ belongs to this $\mathrm{SL}_2$. It fixes $w_1$, but $\phi_1(u)=-1$. Conjugating $T_1$ by $S$ therefore gives centre $w_1$ and functional $\phi_1+\psi$. Multiplying by $T_1^{-1}$ gives $I+w_1\psi$. The independent functionals $\phi_1,\psi$ span $\operatorname{Ann}(w_1)$, so every transvection of centre $w_1$ occurs.

Likewise $T_1,T_2$ generate $\mathrm{SL}$ on $\operatorname{span}(w_1,w_2)$, and $T_1,T_3$ on $\operatorname{span}(w_1,w_3)$, since their cross pairings are nonzero. They send $w_1$ to nonzero multiples of $w_2,w_3$. Conjugation supplies all transvections of these centres as well. Since the centres form a basis, the group contains $\mathrm{SL}_3(\mathbf F_3)$.

Proposition~\ref{prop:determinant-reflections} gives the full determinant $\mathbf F_p^*$, with its cubic-boundary argument in Section~\ref{sec:cubic-generic}. A subgroup containing $\mathrm{SL}_g(\mathbf F_p)$ and having full determinant is $\mathrm{GL}_g(\mathbf F_p)$. The arithmetic image contains this geometric image and lies in the same ambient linear group, so it is equal to it. This proves the odd-characteristic part of Theorem~\ref{thm:plane-images}.

\section{Even plane degree in characteristic two}\label{sec:plane-even}
This section applies the full node-translation subgroup when its node number is odd.

If $p=2$ and $d\ge4$ is even, then $N=d-1$ is odd. The moving-line boundary gives every $I+\operatorname{Hom}(V/P_i,P_i)$. Lemma~\ref{lem:plane-stars} yields
$\mathrm{SL}_g(\mathbf F_2)=\mathrm{GL}_g(\mathbf F_2)$. The arithmetic image is the same by containment. This proves every even degree, not only the smallest generic-matrix case.

\section{Odd plane degree in characteristic two}\label{sec:plane-odd-two}
This section proves that the invariant covector is the entire restriction on the monodromy image.

Suppose $d\ge5$ is odd. Let $\lambda$ be the globally defined nonzero covector from the characteristic-two matrix calculation, and put $K=\ker\lambda$. Both geometric and arithmetic monodromy lie in its stabilizer. At the ordinary line arrangement, $\lambda$ is the sum-of-edge-coefficients functional. Indeed $S_d$ is transitive on edges, so the only invariant line in the dual edge space is generated by the all-ones vector. This vector belongs to graph cycle space exactly when every vertex has even degree. Here the degree is $d-1$, which is even. Hence $V^*$ has exactly this one-dimensional $S_d$-fixed subspace. The nonzero global $\lambda$ is $S_d$-invariant, so it is that vector.

For every vertex, the node sum on $P_i$ is the restriction of $\lambda$. Thus the boundary gives
\[
 N_i=I+\operatorname{Hom}(V/P_i,K_i),\qquad K_i=P_i\cap K,
                      \qquad \dim K_i=d-3.
\]
The spaces $K_i,K_j$ have zero intersection for $i\ne j$: by Lemma~\ref{lem:plane-stars} their possible common edge is $e_{ij}$, and $\lambda(e_{ij})=1$. They span $K$. To see this, the kernel of edge-sum in the edge space is generated by sums of two adjacent edges, because the line graph of $K_d$ is connected. Each such sum lies in a $K_i$. Coboundaries have even edge-sum since $d-1$ is even, so quotienting gives precisely $K$.

\paragraph{The action on the kernel.}
Because $\lambda$ is nonzero on $P_i$, one has $V=K+P_i$ and
\[
 \operatorname{Ann}_V(P_i)\xrightarrow{\sim}\operatorname{Ann}_K(K_i).
\]
Thus $N_i$ acts on $K$ as $I+\operatorname{Hom}(K/K_i,K_i)$.

We use the following elementary observation. Let nonzero subspaces $A_i$ span a vector space $W$, and suppose each $A_i$ has zero intersection with some $A_j$. The groups $I+\operatorname{Hom}(W/A_i,A_i)$ generate $\mathrm{SL}(W)$. Fix nonzero $v\in A=A_i$, and choose $B=A_j$ disjoint from $A$. Pick nonzero $w\in B$ and $\alpha\in\operatorname{Ann}(A)$ with $\alpha(w)=1$. For $\beta\in\operatorname{Ann}(B+\mathbf F_pv)$, direct multiplication gives
\[
 [I+v\alpha,I+w\beta]=I+v\beta.
\]
Here $\beta(v)=0$, and all the other products vanish as required. The available centre-$v$ functionals therefore contain both $\operatorname{Ann}(A)$ and $\operatorname{Ann}(B+\mathbf F_pv)$. Their sum is $\operatorname{Ann}(v)$ because $A\cap(B+\mathbf F_pv)=\mathbf F_pv$. Thus every centre-$v$ transvection occurs. A basis from the union of the $A_i$ gives all elementary matrices. Applied to the pairwise disjoint $K_i$, this proves that the restriction of geometric monodromy to $K$ is $\mathrm{GL}(K)$ over $\mathbf F_2$.

\paragraph{The translation kernel.}
Choosing a vector with $\lambda=1$ identifies $\operatorname{Stab}(\lambda)$ with $K\rtimes\mathrm{GL}(K)$. The kernel of the geometric image $G\to\mathrm{GL}(K)$ is a $\mathrm{GL}(K)$-invariant subgroup of $K$, hence is either zero or all of $K$. Suppose it is zero. Then $G$ is a complement and determines a cocycle
\[
 c:\mathrm{GL}_m(\mathbf F_2)\longrightarrow\mathbf F_2^m,
                    \qquad m=\dim K=g-1\ge5.
\]
Every such cocycle is a coboundary for $m\ge4$, as follows directly from elementary matrices. Write $t_{ij}=I+E_{ij}$. The relation $t_{ij}^2=1$ gives $c(t_{ij})_j=0$. If $l$ differs from $i,j$, choose $k$ different from all three. Commutativity of $t_{ij},t_{kl}$ gives
\[
 e_i c(t_{kl})_j=e_k c(t_{ij})_l,
\]
so $c(t_{ij})_l=0$. Therefore $c(t_{ij})=a_{ij}e_i$. The relation $[t_{ij},t_{jk}]=t_{ik}$ gives $a_{ik}=a_{jk}$: its affine action changes coordinate $i$ by $x_k+a_{jk}$. Thus $a_{ik}$ depends only on $k$. Choose $v$ with those coordinates. Then $c(t_{ij})=(t_{ij}-I)v$ for every elementary transvection. They generate $\mathrm{GL}_m(\mathbf F_2)$, so $c(g)=(g-I)v$ for every $g$. The affine action has a fixed point.

Consequently a complement would fix a vector outside $K$ in $V$. But the fixed space of $N_i$ is exactly $P_i$: the centre space $K_i$ is nonzero and every functional in $\operatorname{Ann}(P_i)$ occurs. Their common fixed space is $\bigcap_iP_i=0$, by Lemma~\ref{lem:plane-stars} and the noncollinearity of distinct edge classes. This contradiction proves that the translation kernel is all of $K$. Hence
\[
 G=\operatorname{Stab}(\lambda)
   \simeq\mathbf F_2^{g-1}\rtimes\mathrm{GL}_{g-1}(\mathbf F_2).
\]
The arithmetic image also preserves the globally defined $\lambda$ and contains this geometric image, so it is the same group. This completes Theorem~\ref{thm:plane-images}.

\section{Positioned homologies in characteristic at least five}\label{sec:higher-homologies}
This section locates the scalar elements of the cubic boundary in one arrangement fibre; their precise positions will allow the full higher-dimensional group to be generated.

Assume throughout Sections~\ref{sec:higher-homologies}--\ref{sec:higher-amalgamation} that $p\ge5$, $n\ge3$ and $d\ge n+2$, and put $m=d-3$. The exact sequences of Section~\ref{sec:cubic-cech} give an actual extension of finite local systems
\begin{equation}\label{eq:higher-extension}
0\longrightarrow P\longrightarrow V_{\mathrm{bdry}}
 \longrightarrow E=\bigoplus_{|I|=n-2}E_I\longrightarrow0.
\end{equation}
Here $E_I=H^1_{\mathrm{et}}(Q\cap\Lambda_I,\mathbf F_p)$. Since $3$ is invertible, each three-point permutation module splits canonically into its constant line and its sum-zero plane. Every incoming differential of the degree-zero intersection complex lands in the constant lines. Thus the second filtration sequence in Section~\ref{sec:cubic-cech-quotients} splits naturally as
\begin{equation}\label{eq:higher-P}
P=P_{\mathrm{const}}\oplus\bigoplus_{|J|=n-1}Q_J,
\qquad Q_J=\mathbf F_p^{Z_J}/\mathbf F_p\mathbf1.
\end{equation}
The constant part has dimension $\binom mn$. This is a splitting of the actual local system, not just a decomposition of its semisimplification.

These sequences extend to the parameter open allowing the cubic itself to be a transverse union of three hyperplanes. Indeed its positive-dimensional strata remain connected; cubic strata of dimension at least two have no positive coherent cohomology by the hypersurface sequence, while its plane sections are ordinary nodal triangles. Its zero-dimensional strata remain three distinct points. The same intersection-complex calculation therefore gives locally free cohomology of the same ranks and invertible Frobenius. Section~\ref{sec:cubic-generic} supplies a finite \'etale Lang system throughout this open. Restriction to the smooth-cubic open leaves its finite image unchanged, by the connected-cover argument there.

Let $H$ be this boundary image. Theorem~\ref{thm:joint} and \eqref{eq:higher-P} give
\begin{equation}\label{eq:higher-group-extension}
1\longrightarrow N\longrightarrow H\longrightarrow
 A=\prod_J S_3\times\prod_I\mathbf F_p^\times\longrightarrow1,
\qquad N\subset\operatorname{Hom}(E,P).
\end{equation}
The kernel acts identically on $P$ and $E$, so it is an additive $\mathbf F_p$-subspace of the displayed Hom space. The order of $A$ is prime to $p$.

There is also a distinguished subgroup $\Gamma\simeq S_3$ in $H$: at a cubic consisting of three general hyperplanes, permute those three unlabelled factors and leave all other factors labelled. The connected ordered configuration cover gives this action on the arrangement fibre. Its action on each $E_I$ is sign, since the cohomology of a triangle is its oriented cycle line. It acts on each $Q_J$ as the standard two-dimensional representation and fixes $P_{\mathrm{const}}$. In particular $P$ has no sign constituent for $\Gamma$.

\subsection{A complement containing the arrangement permutations}
The group extension \eqref{eq:higher-group-extension} has a complement containing $\Gamma$. Here is an elementary verification, including this last qualification. A set-theoretic section has a $2$-cocycle $c(a,b)\in N$. Set
\[
t(a)=|A|^{-1}\sum_{b\in A}c(a,b).
\]
Summing the cocycle identity in the third argument gives
\[
c(a,b)=a\,t(b)-t(ab)+t(a).
\]
Changing the section by $-t$ makes it a group homomorphism. The image of $\Gamma$ in $A$ is injective. The difference between its given lift and its lift in this complement is a $1$-cocycle on $\Gamma$ with coefficients in $N$. Averaging this cocycle expresses it as a coboundary, because $6$ is invertible. Conjugating the complement by the corresponding element of $N$ makes it contain $\Gamma$.

The scalar subgroup on each individual $E_I$ commutes with the image of $\Gamma$ in $A$. Its lift in this complement therefore commutes with $\Gamma$. The $\Gamma$-sign isotypic subspace of $V_{\mathrm{bdry}}$ maps isomorphically onto $E$ in \eqref{eq:higher-extension}, because $P$ has no sign constituent. The lifted scalar acts as the chosen scalar on that one line, as the identity on the other lines, and as the identity on $P$. This identifies its position, beyond the conjugacy class obtained in Proposition~\ref{prop:determinant-reflections}.

\subsection{Exterior coordinates at an arrangement}
Let $E_0=\mathbf F_p^d$ have basis $e_1,\ldots,e_d$, put $u=\sum_i e_i$, and set $W=E_0/\mathbf F_pu$. At a transverse arrangement of $d$ hyperplanes, the incidence resolution gives the $S_d$-equivariant marking
\begin{equation}\label{eq:higher-wedge}
V=\bigwedge^nW.
\end{equation}
Indeed its last differential is $u\wedge-:\bigwedge^{n-1}E_0\to\bigwedge^nE_0$, and there is no next term. The quotient is \eqref{eq:higher-wedge}, of dimension $\binom{d-1}n$. All cohomology is the stable constant incidence part, so this is an ordinary fibre of the same Lang system.

Fix a triple $T=\{a,b,c\}$ and define
\[
A_T=\langle e_a-e_c,e_b-e_c\rangle,
\qquad B_T=\langle e_i:i\notin T\rangle.
\]
Since $3$ is invertible, $W=A_T\oplus B_T$. Permutations of $T$ act standardly on $A_T$ and trivially on $B_T$. Let $S_T$ be their sign isotypic subspace in $V$, and let $P_T$ be the kernel of the sign projector. Then
\begin{equation}\label{eq:higher-block}
S_T=\bigwedge^2A_T\otimes\bigwedge^{n-2}B_T,
\qquad V=S_T\oplus P_T.
\end{equation}
For $I\subset\{1,\ldots,d\}\setminus T$ of size $n-2$, put
\[
v_{T,I}=(e_a-e_c)\wedge(e_b-e_c)\wedge e_I.
\]
These form a basis of $S_T$. Refining the cubic into its three components identifies their images in the elliptic quotient in \eqref{eq:higher-extension}: the cochain supported on $I$ and two of $a,b,c$ is the alternating triangle cochain. Its projection to $E_J$ is zero for $J\ne I$ and is three times an oriented triangle generator for $J=I$. One can read this directly in the double intersection complex: the degree-one refinement term over the plane indexed by $I$ is its triangle complex, and no other outside index set occurs in that term. Since $3$ is invertible, this gives exactly the indicated coordinate lines.

Let $\phi_{T,I}$ be the coordinate dual of $v_{T,I}$ on $S_T$, extended by zero on $P_T$. The scalar construction above supplies the exact elements
\begin{equation}\label{eq:higher-homology}
R_{T,I}(z)=1+(z-1)v_{T,I}\phi_{T,I},
\qquad z\in\mathbf F_p^\times,
\end{equation}
in the geometric image. Conjugation by the arrangement $S_d$ gives them for every triple $T$.

\section{Elementary matrices on the alternating blocks}\label{sec:higher-block-generation}
This section uses explicit overlap coefficients to turn the positioned homologies into a full linear group on each $S_T$.

Fix $T=\{a,b,c\}$, write $k=n-2$, and choose adjacent $k$-subsets of its complement,
\[
I=K\cup\{i\},\qquad L=K\cup\{j\},\qquad |K|=k-1.
\]
Use compatible orientations, placing $e_i,e_j$ before $e_K$. Set
\[
\begin{gathered}
v=v_{T,I},\quad v'=v_{T,L},\quad
T'=\{a,b,j\},\quad J=K\cup\{c\},\quad w=v_{T',J},\\
\phi=\phi_{T,I},\quad\phi'=\phi_{T,L},\quad\psi=\phi_{T',J}.
\end{gathered}
\]
The required coefficients are
\begin{equation}\label{eq:higher-overlap}
\phi(w)=-\frac13,\qquad\psi(v)=0,
\qquad\psi(v')=-1,\qquad\phi(v')=0.
\end{equation}
For a complete calculation, put $A_0=e_a-e_c$, $B_0=e_b-e_c$ and $s=\sum_{r\notin T}e_r$. Then $e_c=-(A_0+B_0+s)/3$. The projection onto $S_T$ of
\[
((e_a-e_j)\wedge(e_b-e_j))\wedge e_c\wedge e_K
\]
is
\[
-\frac13 A_0\wedge B_0\wedge s\wedge e_K
-\frac23 A_0\wedge B_0\wedge e_j\wedge e_K.
\]
Its coefficients at $v,v'$ are $-1/3,-1$. Interchanging $c,j$ in the same calculation gives $\psi(v')=-1$. Finally every row of the coordinate determinant defining $\psi$ vanishes on $e_i$, since $i\notin T'\cup J$, so $\psi(v)=0$. This proves \eqref{eq:higher-overlap}, with any additional outside indices present.

Choose $a_0,b_0\ne0,1$. Conjugate $R_w(b_0)$ by $R_v(a_0)$ and multiply by $R_w(b_0)^{-1}$. Since $\psi(v)=0$, the result is
\[
1+(1-b_0^{-1})(a_0-1)\phi(w)\,v\psi.
\]
Its coefficient is nonzero, and powers give $1+t v\psi$ for every $t\in\mathbf F_p$. Conjugating by $R_{v'}(c_0)$, with $c_0\ne0,1$, fixes the centre $v$ and changes its row by
\[
(c_0^{-1}-1)\psi(v')\phi'.
\]
Multiplication by the inverse original transvection gives a nonzero multiple of $v\phi'$. Powers again supply
\begin{equation}\label{eq:higher-elementary}
1+t v_{T,I}\phi_{T,L}\qquad(t\in\mathbf F_p).
\end{equation}
The graph of $k$-subsets joined by a one-index replacement is connected. Commutators of consecutive elementary matrices along a simple path give every ordered pair of distinct basis indices. Together with the diagonal scalings \eqref{eq:higher-homology}, they generate
\begin{equation}\label{eq:higher-full-block}
G_T=\operatorname{GL}(S_T)\quad\hbox{acting identically on }P_T.
\end{equation}
Here $1\le k\le m-1$, so every block has dimension at least two.

\section{Amalgamating the blocks}\label{sec:higher-amalgamation}
This section proves an elementary generation lemma and verifies its hypotheses for all higher-dimensional arrangement blocks, including when $p$ divides $d$.

\begin{lemma}\label{lem:higher-amalgamation}
Let $p$ be odd and let $V=S_i\oplus P_i$ be a collection of decompositions over $\mathbf F_p$, with $\dim S_i\ge2$. Put $L_i=\operatorname{Ann}(P_i)\subset V^*$. Suppose the graph joining $i,j$ when $S_i\cap S_j\ne0$ is connected, the $S_i$ span $V$, and the $L_i$ span $V^*$. Then the groups $\operatorname{GL}(S_i)$ acting identically on $P_i$ generate $\operatorname{GL}(V)$.
\end{lemma}
\begin{proof}
For $v\ne0$ and a row space $L$ evaluating nontrivially on $v$, define
\[
\mathcal A(v,L)=\{1+v\phi:\phi\in L,\ 1+\phi(v)\ne0\}.
\]
If $L,M$ both evaluate nontrivially on $v$, the two axis groups generate $\mathcal A(v,L+M)$. To prove this, choose $\lambda\in L$, $\mu\in M$ with $\lambda(v)=\mu(v)=1$. All unipotent rows in $L\cap\operatorname{Ann}(v)$ and $M\cap\operatorname{Ann}(v)$ are present. For $a\ne0,1$,
\[
[1+(a-1)v\lambda][1+(a-1)v\mu]^{-1}
 =1+(a-1)v(\lambda-\mu).
\]
Powers give every $\mathbf F_p$ multiple of this row difference. These three spaces span $(L+M)\cap\operatorname{Ann}(v)$, and the original axis scalars supply every nonzero value of $1+\phi(v)$. This proves the assertion.

Start with one block. Its group supplies $\mathcal A(v,L_i)$ for every nonzero $v\in S_i$. Inductively suppose this holds on every block of a connected collection with the full row space $L_C=\sum_{i\in C}L_i$. Add a block $S_j$ meeting $S_i$ in a nonzero $v$. The axis assertion first gives $\mathcal A(v,L_C+L_j)$. Each old block group preserves $L_C+L_j$: if $g$ acts identically on $P_i$, then $\phi g-\phi\in L_i$. Transitivity of $\operatorname{GL}(S_i)$ on its nonzero vectors supplies the enlarged axis group for every centre in $S_i$. Propagate along the old connected intersection graph, and use the new block group to obtain it on $S_j$ too. Induction gives axis rows from the total sum $V^*$ at every nonzero centre in every block.

Choose a basis of $V$ from the union of the $S_i$. All elementary transvections and all individual diagonal scalings in that basis are now present. They generate $\operatorname{GL}(V)$.
\end{proof}

We verify the three hypotheses for \eqref{eq:higher-full-block}. Intrinsically,
\[
S_T=(\bigwedge^2 A_T)\wedge\bigwedge^{n-2}W.
\]
If $T,T'$ share two vertices, $A_T+A_{T'}$ has dimension three. The nonzero space
\[
\bigwedge^3(A_T+A_{T'})\wedge\bigwedge^{n-3}W
\]
lies in $S_T\cap S_{T'}$. The graph of triples joined by a two-vertex overlap is connected. This proves the intersection hypothesis.

Let $D\subset W$ be spanned by all differences $e_i-e_j$. Its codimension is at most one. The spaces $\bigwedge^2 A_T$ span $\bigwedge^2D$, as is seen from the triples $\{i,j,d\}$. Hence the $S_T$ span
\[
(\bigwedge^2D)\wedge\bigwedge^{n-2}W.
\]
Every basis wedge of degree $n\ge3$ contains at least $n-1\ge2$ factors from a codimension-at-most-one subspace. The displayed span is therefore all $\bigwedge^nW$. This argument includes $p\mid d$, when $D$ has codimension one.

For the dual assertion, put $r=d-1$ and $q=r-n$. Exterior duality gives
\[
V^*\simeq\bigwedge^qW\otimes(\det W)^*.
\]
On the permutations of $T$, the determinant of $W$ is sign. Thus the sign isotypic space $L_T=\operatorname{Ann}(P_T)$ corresponds to the invariant subspace in $\bigwedge^qW$, tensored with $(\det W)^*$. Because $W=A_T\oplus B_T$ and $p>3$, that invariant subspace is exactly $\bigwedge^qB_T$: $A_T$ has no invariants and $\bigwedge^2A_T$ is sign. The spaces $\bigwedge^qB_T$ span $\bigwedge^qW$. Indeed $q=d-1-n\le d-4$, and every wedge of $q$ coordinate classes lies outside some triple $T$. Those coordinate wedges span the whole exterior power.

Lemma~\ref{lem:higher-amalgamation} now gives the full geometric $\operatorname{GL}_{\binom{d-1}n}(\mathbf F_p)$. The arithmetic image contains it and lies in the same group, so is equal. This proves Theorem~\ref{thm:higher-images} for $p\ge5$.


\section{A generic-matrix chart in characteristic three}\label{sec:three-seed}
This section proves full monodromy in characteristic three at $d=n+2$ by making the entire Hasse--Witt matrix a set of affine coordinates on an explicit coefficient open.

Assume $n\ge3$, $d=n+2$ and $p=3$. The positive indices are $a_i=\mathbf1+e_i$, $0\le i\le n$, and $h=n+1$. The convention of \eqref{eq:pn-1-4} gives
\begin{equation}\label{eq:three-matrix}
B_3(f)_{i,j}=[x^{M_{i,j}}]f^2,
\qquad M_{i,j}=2\mathbf1+3e_j-e_i.
\end{equation}
For $i=j$, the target exponent has entry $4$ at $j$ and entry $2$ elsewhere. For $i\ne j$, it has entry $5$ at $j$, entry $1$ at $i$, and entry $2$ elsewhere. In particular it has at most one entry at least four.

Call a degree-$d$ monomial \emph{low} if every exponent is at most three, and \emph{high} otherwise. Write $f=g+z$ accordingly. Two high monomials have either an exponent sum at least eight in one coordinate, or exponents at least four in two coordinates. Their product can never have exponent $M_{i,j}$. Consequently
\begin{equation}\label{eq:three-linear}
B_3(g+z)=B_3(g)+2L_g(z),
\qquad L_g(z)_{i,j}=[x^{M_{i,j}}]gz.
\end{equation}
A high monomial with exponent at least four in position $j$ can contribute only to column $j$. A monomial with two such positions contributes to no column.

\subsection{Explicit invertible coefficient minors}
Fix $j$ and choose three distinct other indices $a,b,c$. Let $R$ be the remaining set, of size $n-3$, and put $P_R=\prod_{r\in R}x_r$. Use the low polynomial
\[
g_j=(x_ax_b^2x_c^2+x_jx_a^2x_c^2+x_jx_a^2x_b^2)P_R.
\]
It has degree $n+2$ and all its exponents are at most two. Select the following $h$ high monomials:
\begin{equation}\label{eq:three-monomials}
\begin{aligned}
Z_0&=x_j^5P_R,&Z_a&=x_j^4x_aP_R,\\
Z_b&=x_j^4x_bP_R,&Z_c&=x_j^4x_cP_R,\\
Z_r&=x_j^5x_a\prod_{s\in R\setminus\{r\}}x_s
 &&(r\in R).
\end{aligned}
\end{equation}
They are distinct degree-$n+2$ monomials, each with precisely one exponent at least four, in position $j$. Ordering rows as $j,a,b,c,R$ and columns as $0,a,b,c,R$, the selected matrix of $L_{g_j}$ in column $j$ is
\begin{equation}\label{eq:three-minor}
\begin{pmatrix}
0&1&0&0\\1&0&0&0\\0&0&1&0\\0&0&0&1
\end{pmatrix}\oplus I_{n-3}.
\end{equation}
Write $\nu(Z)$ for the exponent vector of a monomial $Z$. Here is the complete coefficient check. For the first four rows and columns, each variable in $R$ has target exponent two and selected exponent one, so the remaining exponent is one. The calculation reduces to the four variables $j,a,b,c$. Its four nonzero differences are
\[
\begin{aligned}
M_{j,j}-\nu Z_a&=\nu(x_ax_b^2x_c^2P_R),&
M_{a,j}-\nu Z_0&=\nu(x_ax_b^2x_c^2P_R),\\
M_{b,j}-\nu Z_b&=\nu(x_jx_a^2x_c^2P_R),&
M_{c,j}-\nu Z_c&=\nu(x_jx_a^2x_b^2P_R),
\end{aligned}
\]
Every other difference has a negative exponent or is absent from $g_j$. For row $r\in R$ and column $Z_r$, the difference is the first monomial of $g_j$. For a different row and column $Z_r$, coordinate $r$ would require exponent two in $g_j$, whereas its exponent there is one. Conversely row $r$ and a first-four column would require exponent zero there. This proves \eqref{eq:three-minor}; its determinant is $-1$.

For the general low polynomial, let $D_j$ be this minor determinant. It is a nonzero polynomial. The product $\prod_jD_j$ is nonzero in the integral low-coefficient polynomial ring, so there is a nonempty open $B_{\mathrm{low}}$ on which every column minor is invertible. The special polynomials $g_j$ used to verify the separate minors need not be the same polynomial.

\subsection{The coefficient chart and its connected Lang cover}
The chosen monomials are distinct between columns, because their unique exponent at least four occurs at different positions. Use their $h^2$ coefficients as distinguished high coefficients and put the remaining high coefficients in an affine space $B_{\mathrm{high}}$. Set $B=B_{\mathrm{low}}\times B_{\mathrm{high}}$.

Equation~\eqref{eq:three-linear} gives an isomorphism
\begin{equation}\label{eq:three-chart}
B\times\mathbf A^{h^2}\xrightarrow{\ \sim\ }B\times\operatorname{Mat}_h,
\qquad f\longmapsto (\text{remaining coefficients},B_3(f)).
\end{equation}
Its linear part in the distinguished coefficients is block diagonal by columns, with blocks twice the invertible minors $D_j$. The inverse solves those linear systems, using only the determinants already inverted. Other high coefficients only change the known affine constant term. Thus \eqref{eq:three-chart} is a regular coordinate isomorphism, not just a dominant map. It describes an open of the full affine coefficient space; no coefficient has been fixed to zero.

The ordinary part is $B\times\operatorname{GL}_h$. Its Lang basis cover, by the elementary fixed-vector calculation of Section~\ref{sec:pn-2}, has total space $B\times\operatorname{GL}_h$, with map
\[
(b,T)\longmapsto\bigl(b,T(T^{[3]})^{-1}\bigr).
\]
It is a finite \'etale $\operatorname{GL}_h(\mathbf F_3)$-torsor and its total space is geometrically integral. Hence its geometric image is the full group.

The ordinary coefficient space is integral and normal. Its smooth hypersurface open is nonempty by Section~\ref{sec:pn-3}, and it meets the nonempty coefficient chart above. Restricting the connected finite \'etale basis cover to this nonempty open preserves connectedness. Passing from nonzero affine equations to projective equations places this full affine image inside the projective image. The latter is therefore also $\operatorname{GL}_h(\mathbf F_3)$. The arithmetic image contains it and is bounded by the same group. This proves the degree-$d=n+2$ characteristic-three case of Theorem~\ref{thm:higher-images}.


\section{The modular cubic boundary in characteristic three}\label{sec:higher-three-modular}
This section identifies the nonsplit node module that forces the additional unipotent transformations in characteristic three. Assume $n\ge3$, $d\ge n+3$, and put $m=d-3$.

The cubic-boundary calculation of Section~\ref{sec:cubic-cech} gives
\begin{equation}\label{eq:three-extension}
0\longrightarrow P\longrightarrow V\longrightarrow
E=\bigoplus_{|I|=n-2}E_I\longrightarrow0,
\qquad
0\longrightarrow P_0\longrightarrow P\longrightarrow
\bigoplus_{|J|=n-1}Q_J\longrightarrow0.
\end{equation}
Here $P_0$ is constant of dimension $\binom mn$ and $Q_J=\mathbf F_3[Z_J]/\mathbf F_3\mathbf1$. We do not split the second sequence. In characteristic three its failure to split is essential.

For clarity, the first sequence follows directly from the intersection resolution. A cubic of dimension at least two has no positive structure-sheaf cohomology; a plane cubic has one-dimensional $H^1$; and a line section has length three. These statements follow from the hypersurface sequence and the projective-space calculation in \citet[Lemmas~30.8.1--30.8.2]{StacksProjective}. Thus the only positive row is the elliptic row at bidegree $(n-2,1)$. The resolution ends in degree $n-1$, so no differential can leave that row toward degree $n$ in the zero row. Replacing the three point values by one diagonal constant gives the constant skeleton complex, whose top cohomology is $P_0$. The quotient is supported only in the last degree, giving the second sequence in~\eqref{eq:three-extension}. In particular the entire action on $P$ is determined by actual permutations of the three-point sets, since all other zero-row terms and maps are constant.

Allow the cubic to become three labelled hyperplanes $a,b,c$, transverse to the other $m$ hyperplanes. The plane sections are ordinary triangles. The grouped intersection resolution remains exact: locally its cubic equation is a product of coordinate variables disjoint from those defining the remaining hyperplanes, and the alternating ideal complex is exact monomial by monomial. Refining the grouped resolution gives the fully labelled one. The hypersurface calculation keeps all ranks constant across this specialization and makes the cohomology and intersection maps commute with base change. The ordinary coefficient open containing the triangle and the smooth-cubic open is normal and integral. Section~\ref{sec:cubic-generic} therefore identifies the finite image at this common ordinary fibre.

Let
\[
U_0=\mathbf F_3^d,\qquad u=\sum_l e_l,\qquad
W=U_0/\mathbf F_3u.
\]
The labelled arrangement complex has differential $u\wedge$, and hence
\begin{equation}\label{eq:three-arrangement}
V=\bigwedge^nU_0/(u\wedge\bigwedge^{n-1}U_0)=\bigwedge^nW.
\end{equation}
An $n$-fold point cochain maps to the exterior product of its labels. The parameter space of ordered general-position hyperplanes is irreducible; quotienting the three labels $a,b,c$ gives a connected $S_3$ cover. Their physical permutations are therefore actual elements of the grouped boundary image, acting on~\eqref{eq:three-arrangement} by exterior powers.

Put $A=e_a$, $B=e_b$, let $B_0$ be spanned by the $m$ outside vectors, and put $S=\sum_{j=1}^m e_j$. Then $e_c=-A-B-S$ and $W=\mathbf F_3A\oplus\mathbf F_3B\oplus B_0$. Refinement sends the three values at $Z_J$ to $A\wedge e_J,B\wedge e_J,e_c\wedge e_J$ and the hyperplane-only points to their outside cochains. Thus
\begin{align}
P={}&\bigwedge^nB_0\oplus
 A\wedge\bigwedge^{n-1}B_0\oplus
 B\wedge\bigwedge^{n-1}B_0,\label{eq:three-P}\\
P_0={}&\bigwedge^nB_0,
\qquad E=\bigoplus_I\mathbf F_3\overline w_I,
\qquad w_I=A\wedge B\wedge e_I.\nonumber
\end{align}
The outside $n$-fold cochains together with $A\wedge e_J$ and $B\wedge e_J$ form the displayed independent basis; the $e_c$ values satisfy $e_c=-A-B-S$. Its dimension agrees with~\eqref{eq:three-extension}, so these identifications are exact. Let $\phi_I$ denote the coordinate row of $V\to E_I$ in this quotient basis.

At an individual node index $J$, the transposition $(a\ b)$ swaps $A\wedge e_J$ and $B\wedge e_J$, fixes all other such pairs, and fixes $P_0$. The action on each $Q_J$ is faithful. Theorem~\ref{thm:joint} therefore identifies the quotient of the actual boundary image acting on $P$ and $E$ with
\begin{equation}\label{eq:three-graded-group}
L=\prod_J S_3\times\prod_I\{1,-1\}.
\end{equation}
This specifies the actual modular node action, not just its composition factors.

\section{Separating the elliptic characters in characteristic three}\label{sec:higher-three-kernel}
This section splits the finite group between $P$ and $E$, while preserving the nonsplit node module inside $P$.

Let $G_B$ be the boundary image and let $N$ be the kernel of its action on $P$ and $E$. Its elements have the form $1+b$ with $b:E\to P$, so $N$ is an additive $\mathbf F_3$ subspace of $\operatorname{Hom}(E,P)$. Lift the element of~\eqref{eq:three-graded-group} acting trivially on $P$ and by $-1$ on every $E_I$. Every such lift has block form $\left(\begin{smallmatrix}1&b\\0&-1\end{smallmatrix}\right)$ and square one. Use its negative eigenspace as a complement to $P$, so that the lift $c_0$ is $\operatorname{diag}(1,-1)$.

For $x=\left(\begin{smallmatrix}A_0&b\\0&D\end{smallmatrix}\right)\in G_B$, direct multiplication gives
\[
[x,c_0]=1+2bD^{-1}\in N.
\]
Thus $bD^{-1}\in N$, and multiplying $x$ on the left by $1-bD^{-1}$ removes its off-diagonal block. The centralizer of $c_0$ maps onto $L$ and meets $N$ trivially. Consequently $G_B$ contains a block-diagonal copy of $L$, and $N$ is stable under $b\mapsto A_0bD^{-1}$.

The independent sign on $E_I$ projects $N$ to its $I$-th column: subtract the original element from its conjugate, which multiplies that column by $-2=1$. Put
\[
N_I=\{v\in P:v\phi_I\in N\}.
\]
Each $N_I$ is stable under every independent node permutation. None of this asserts a splitting of $P_0\subset P$.

\section{Modular norms force actual transvections}\label{sec:higher-three-force}
This section obtains supported rank-one transformations from a physical three-cycle and identifies their rows as arrangement cochains.

Let $\sigma$ be the physical cycle $a\mapsto b\mapsto c\mapsto a$. It acts by $+1$ on every triangle quotient $E_I$. In the chosen complement, write its lift of $\overline w_I$ as $x_I=w_I+p_I$, $p_I\in P$. The off-diagonal column of $\sigma$ is $b_I=(\sigma-1)x_I$. Section~\ref{sec:higher-three-kernel} gives $b_I\in N_I$, and stability gives
\[
z_I=(\sigma-1)b_I=(\sigma-1)^2x_I=(1+\sigma+\sigma^2)x_I\in N_I.
\]
The norm sends $P$ into $P_0$: it is zero on $P_0$ and sends $A\wedge e_J$ and $B\wedge e_J$ to $-S\wedge e_J$. Moreover
\[
\sigma w_I=w_I-B\wedge S\wedge e_I,
\qquad \sigma^2w_I=w_I+A\wedge S\wedge e_I.
\]
Therefore
\begin{equation}\label{eq:three-norm}
z_I=(A-B)\wedge S\wedge e_I+q_I,
\qquad q_I\in P_0.
\end{equation}
For $j\notin I$ and $J=I\cup\{j\}$, apply $t_J-1$, where $t_J$ is the independent node transposition $(a\ b)$ at $J$. It kills all other node components and $P_0$, and multiplies the displayed $J$ component by $-2=1$. Hence $(A-B)\wedge e_J\in N_I$. The actual image contains
\begin{equation}\label{eq:three-supported}
1+\bigl[(e_a-e_b)\wedge e_I\wedge e_j\bigr]\phi_I.
\end{equation}
Its centre lies in $P$ and its row vanishes there, so this is a transvection on the full representation. Regrouping any three of the labelled hyperplanes as the cubic gives the same conclusion for every allowed choice of labels at the fixed fibre~\eqref{eq:three-arrangement}.

For an ordered $(n+1)$-element label set $U$, let $\delta_U$ have the alternating boundary values on its $n$-faces and value zero on the other coordinate $n$-products. It annihilates $u\wedge e_K$, since the two possible boundary terms cancel. Thus it is a row on $\bigwedge^nW$. For $U=\{a,b,c\}\cup I$, it vanishes on~\eqref{eq:three-P}; on $w_{I'}$ it is zero unless $I'=I$, and then is $1$ or $-1$. Accordingly
\begin{equation}\label{eq:three-row}
\phi_I=\pm\delta_{\{a,b,c\}\cup I}.
\end{equation}
The sign is absorbed by taking the inverse transvection. This quotient-row calculation is independent of the complement used to exhibit $N$.

\section{Elementary generation in characteristic three}\label{sec:higher-three-group}
This section converts the supported transformations into all elementary matrices, completing the characteristic-three image calculation without a classification theorem for finite linear groups.

Choose $e_1,\ldots,e_{d-1}$ as a basis of $W$, so $e_d=-\sum_{r<d}e_r$. Coordinate vectors of $V$ are $e_A$ for $n$-subsets $A$ of the first $d-1$ labels. Fix $R=\{r_0,\ldots,r_n\}$ in increasing order, and write
\[
f_i=(-1)^i e_{R\setminus\{r_i\}}.
\]
For distinct $i,j,k$, apply~\eqref{eq:three-supported} using $a=r_j$, $b=r_k$, $c=d$, outside set $I=R\setminus\{r_i,r_j,r_k\}$, and the additional label $r_i$. Its row in~\eqref{eq:three-row} is a nonzero multiple of $f_i^*$ on the entire coordinate space: only the face $R\setminus\{r_i\}$ avoids $d$. Its centre is a common nonzero multiple of $f_j+f_k$, by the alternating face signs. Taking an inverse if necessary gives every
\begin{equation}\label{eq:three-triple}
1+(f_j+f_k)f_i^*.
\end{equation}
These transformations fix all coordinates outside this clique of $n+1$ faces.

Since $n+1\ge4$, for each $i\ne j$ choose two other distinct indices $k,l$. The same-row transformations commute, and
\begin{align*}
&[1+(f_j+f_k)f_i^*]\,[1+(f_j+f_l)f_i^*]\,[1-(f_k+f_l)f_i^*]\\
&\hspace{35mm}=1+2f_jf_i^*.
\end{align*}
Squaring gives $1+f_jf_i^*$. Thus the image contains all elementary transformations within each clique.

The graph of coordinate $n$-subsets joined by one-element replacements is connected. Every edge belongs to the clique of faces of its $(n+1)$-element union, so both elementary directions occur along every edge. For distinct coordinate vertices $A,B,C$,
\[
[1+E_{CB},1+E_{BA}]=1+E_{CA}.
\]
Commuting along a simple path gives every off-diagonal elementary matrix. They generate $\mathrm{SL}_h(\mathbf F_3)$ by Gaussian elimination; the determinant-one diagonal operations follow from the usual two-by-two elementary identities.

Finally the block-diagonal complement of Section~\ref{sec:higher-three-kernel} contains the element acting by $-1$ on one $E_I$ and identically on $P$ and all other elliptic lines. Its determinant is $-1$. The image therefore contains $\mathrm{SL}_h(\mathbf F_3)$ and both determinants, hence equals $\mathrm{GL}_h(\mathbf F_3)$. The geometric-to-arithmetic containment proves the same equality over every perfect ground field. Together with the first-degree chart in Section~\ref{sec:three-seed}, this completes the characteristic-three case of Theorem~\ref{thm:higher-images}.


\section{A local elliptic cover in characteristic two}\label{sec:higher-two-local}
This section constructs the local translations that will couple an elliptic quotient to its three-point node factors. Throughout Sections~\ref{sec:higher-two-local}--\ref{sec:higher-two-group}, the characteristic is two, $n\ge3$, and $d\ge n+2$.

Write $d=m+3$. The cubic boundary of Sections~\ref{sec:cubic-cech} and~\ref{sec:cubic-generic} has an exact sequence
\begin{equation}\label{eq:two-extension}
0\longrightarrow P\longrightarrow V\longrightarrow
E=\bigoplus_{|I|=n-2} E_I\longrightarrow0,
\qquad E_I=H^1_{\mathrm{et}}(C_I,\mathbf F_2).
\end{equation}
Each $E_I$ is constant of rank one. Since $3$ is invertible, the point-value modules split by the trace, giving
\[
P=P_{\mathrm{const}}\oplus\bigoplus_{|J|=n-1}Q_J,
\qquad Q_J=\mathbf F_2[Z_J]/\mathbf F_2\mathbf1
       =\ker\bigl(\operatorname{sum}:\mathbf F_2[Z_J]\to\mathbf F_2\bigr).
\]
The image on $P$ is the full product of the node groups $S_3$, by Theorem~\ref{thm:joint}. Its kernel on $P$ and $E$ is an additive subgroup of $\operatorname{Hom}(E,P)$. For $I\subset J$, let $\phi_I:V\to E_I=\mathbf F_2$ be the quotient coordinate. We shall prove that the actual image contains
\begin{equation}\label{eq:two-supported}
1+v\phi_I\qquad\text{for every }v\in Q_J.
\end{equation}
The assertion concerns the full extension~\eqref{eq:two-extension}.

Let $K=k((t))$, with $k$ algebraically closed, and consider
\[
C_t:\ y^2+txy+y=x^3.
\]
Its projective equation is $y^2z+txyz+yz^2=x^3$. At $t=0$ its partial derivatives are $x^2,z^2,y^2$, so the special cubic is smooth. The generic cubic is smooth and ordinary: its Hasse coefficient is $t$, whereas the special cubic is supersingular.

\begin{lemma}\label{lem:two-local-cover}
The equation
\begin{equation}\label{eq:two-local-AS}
w^2+w=x/t^2
\end{equation}
defines a geometrically connected \textup{\'etale} double cover of $C_t$. For $b\in k\setminus\{0,1\}$, the three points of $C_t\cap\{y=b\}$ are $K$-rational. The Galois action on their lifts to this cover is the full group
\[
\{(\epsilon_0,\epsilon_1,\epsilon_2)\in\mathbf F_2^3:
\epsilon_0+\epsilon_1+\epsilon_2=0\}.
\]
\end{lemma}
\begin{proof}
Only the point $O$ at infinity can be a pole of the right side. With $u=x/y$ at $O$, the cubic equation gives
\[
x=u^{-2}+tu^{-1}+\frac1{ux}.
\]
Here $x$ has pole order two, so the last term is regular at $O$. Subtracting the Artin--Schreier coboundary of $1/(tu)$ makes $x/t^2$ regular there. The defining polynomial has derivative one, proving that the cover is \textup{\'etale} everywhere. This is the elementary local form of the Artin--Schreier construction in \citet[Section~59.63]{StacksAS}.

If the cover were geometrically split, we would have $x/t^2=h^2+h$ over $\overline K(C_t)$. The function $h$ could have no pole except $O$, and its pole there would be simple. It would define a degree-one map from a smooth genus-one curve to $\mathbf P^1$, which is impossible. The genus-one assertion follows from the plane-cubic hypersurface sequence, as in \citet[Lemma~53.9.3]{StacksPlane}. Constants cannot give the nonconstant right side. The cover is therefore geometrically connected and represents the nonzero geometric \textup{\'etale} class.

Put $c=b^2+b\ne0$. The intersection coordinates satisfy
\[
x^3+tbx+c=0.
\]
At $t=0$ its three simple roots are $u_i=\alpha\zeta^i$, where $\alpha^3=c$ and $\zeta$ has order three. Hensel lifting and differentiation give
\[
x_i(t)=u_i+\frac b{u_i}t+O(t^2).
\]
After subtracting the coboundary of $\sqrt{u_i}/t$, the Artin--Schreier class of the fibre over $x_i(t)$ is
\[
\frac{a_i}{t},\qquad
 a_i=\frac b{u_i}+\sqrt{u_i}
     =\frac{u_i^2(b+\sqrt c)}c.
\]
The regular remainder is a coboundary over $k[[t]]$, by Hensel's lemma. The common coefficient is nonzero because $(b+\sqrt c)^2=b$. Thus the three $a_i$ are $q,q\zeta,q\zeta^2$, in some order, with $q\ne0$. They sum to zero and span a two-dimensional $\mathbf F_2$ space. No nonzero $a/t$ is a coboundary in $K$: a pole of $h^2+h$ has even order. Consequently the three fibre characters have exactly one relation, their sum, and their joint image is the displayed balanced group. Its map to $\mathbf F_2^3/\mathbf F_2\mathbf1$ is an isomorphism, since three is odd.
\end{proof}

\section{Supported shears on the higher-dimensional extension}\label{sec:higher-two-support}
This section transports the explicit curve calculation into the full cubic-boundary cohomology and removes all unwanted node components.

Fix $I$ of size $n-2$, choose $j\notin I$, and put $J=I\cup\{j\}$. Choose the hyperplanes and coordinates so that $\Lambda_I$ is the plane of Lemma~\ref{lem:two-local-cover} and $H_j$ restricts to $y=b$. The remaining hyperplanes can be chosen in general position with all induced lines transverse to $C_0$ and all fixed zero-dimensional strata disjoint from $C_0$.

There is an ambient cubic $Q_0$ restricting to $C_0$ on $\Lambda_I$, smooth and transverse on every relevant stratum, with every other plane section smooth ordinary. Here the space of extensions is an irreducible affine space. Another plane meets $\Lambda_I$ in a line, point, or empty set. In the line case its prescribed binary cubic is squarefree, and Section~\ref{sec:plane-cubic-restriction} supplies a smooth ordinary cubic with exactly that restriction. In the point case choose an ordinary cubic avoiding that point and scale its value; the empty case has no restriction. Compatible forms on two linear spaces extend by Section~\ref{sec:cubic-restrictions}. Thus each other plane-ordinarity condition is a nonempty open condition in the extension space.

Smoothness and transversality impose further nonempty opens. Away from the prescribed plane, variations $l^2L$, where $l$ vanishes on the plane but not at the chosen point and $L$ is linear, span all first jets. The incidence of singular points therefore has codimension at least one after projection. On the prescribed smooth section a tangent derivative is already nonzero. The same argument works on each linear stratum; line transversality and avoidance of the finite point strata are also nonempty opens. Their finite intersection is nonempty. Choose $Q_0$ there and lift the variation $txyz$ to an ambient cubic $tQ_1$.

For the trait $Q_t=Q_0+tQ_1$, every node cover remains \textup{\'etale}, every other elliptic section remains ordinary, and precisely the chosen plane section has Hasse coefficient $t$. Its generic point lies in the ordinary boundary open.

\begin{proposition}\label{prop:two-supported}
For every $I\subset J$ as above, the actual boundary image contains all transformations~\eqref{eq:two-supported}.
\end{proposition}
\begin{proof}
The coherent intersection sequences extend over $k[[t]]$ with constant ranks, by the projective hypersurface cohomology calculation. Let $V_0$ be the preimage of the sum of all elliptic summands other than $E_I$. Frobenius on $V_0$ is invertible even at $t=0$: its filtration consists of the constant and node parts and the other ordinary elliptic lines. Its finite Lang local system therefore extends over the strictly henselian trait and has trivial local Galois action. The generic elliptic quotient is also constant, since each line is over $\mathbf F_2$. Every local element consequently acts on the full $V$ as
\begin{equation}\label{eq:two-local-shear}
1+v\phi_I,\qquad v\in P.
\end{equation}
The centre lies in $P$ because the entire elliptic quotient is fixed.

To identify its projection to $Q_J$, select in the coherent intersection complex all summands whose component index sets contain $I$. They form a subcomplex: the differential only adds indices. On $\Lambda_I$ this is
\[
\mathcal O_{\Lambda_I}\longrightarrow
\bigoplus\mathcal O_{\text{curve components}}\longrightarrow
\bigoplus\mathcal O_{\text{nodes}},
\]
in degrees $n-3,n-2,n-1$, for the nodal curve
\[
Z=C_I\cup\bigcup_{l\notin I}(H_l\cap\Lambda_I).
\]
The augmented normalization complex and $H^1(\mathbf P^2,\mathcal O)=H^2(\mathbf P^2,\mathcal O)=0$ identify its top hypercohomology with $H^1(Z,\mathcal O_Z)$. Inclusion into the full complex gives a Frobenius-compatible map
\[
H^1(Z,\mathcal O_Z)\longrightarrow H^{n-1}(Y,\mathcal O_Y).
\]
It is the inclusion of the chosen elliptic quotient and the identity on the three node values at $J$ modulo constants. Constant component cochains contribute only the diagonal, so it is also the identity on the canonical trace-zero factor $Q_J$.

Restriction to $C_I\cup\{y=b\}$ gives its normalization extension
\[
0\longrightarrow\mathbf F_2^3/\mathbf F_2\mathbf1
\longrightarrow H^1_{\mathrm{et}}(C_I\cup\{y=b\},\mathbf F_2)
\longrightarrow H^1_{\mathrm{et}}(C_I,\mathbf F_2)\longrightarrow0.
\]
A lift of the last generator is obtained by gluing its double-cover torsor to the trivial torsor on the rational line, using chosen lifts of the three nodes. Changing the lifts changes the extension coordinate by their three translations modulo the diagonal, as proved in Section~\ref{sec:plane-transvections}. Lemma~\ref{lem:two-local-cover} therefore makes the projection of the centres in~\eqref{eq:two-local-shear} onto $Q_J$ surjective.

Conjugate a local shear by an element supported on the selected $S_3$ factor on the graded pieces, whose existence follows from Theorem~\ref{thm:joint}. The additive kernel $\operatorname{Hom}(E,P)$ is acted on only through those graded actions. The commutator has centre $(\sigma-1)v$ supported on $Q_J$ and still has row $\phi_I$. The two-dimensional trace-zero three-point module is irreducible over $\mathbf F_2$: its three nonzero vectors are permuted transitively by $S_3$. The differences of a nonzero vector span it. Products of same-row shears add their centres, giving all of~\eqref{eq:two-supported}.
\end{proof}

\section{Marking the characteristic-two shears}\label{sec:higher-two-marking}
This section identifies the supported transformations in one common arrangement basis, including their quotient rows.

Degenerate the cubic to three labelled hyperplanes $a,b,c$, transverse to the remaining hyperplanes. Its plane sections are ordinary triangles and its line sections remain three distinct points. The coherent ranks do not change. Locally the cubic equation is a product of coordinate variables disjoint from those of the remaining hyperplanes. Checking monomials proves exactness of the grouped intersection resolution, and its finite double-complex refinement agrees with the fully labelled resolution. All maps commute with Frobenius. The normal-open Lang comparison of Section~\ref{sec:cubic-generic} therefore permits a common marked ordinary fibre.

Put $W=\mathbf F_2^d/\mathbf F_2\mathbf1$. At this fibre
\begin{equation}\label{eq:two-arrangement}
V=\bigwedge^nW,
\end{equation}
because the last cochain differential is exterior multiplication by $\sum e_i$. Unordered arrangements give the actual $S_d$ permutation action. For $T=\{a,b,c\}$ and $J=I\cup\{j\}$, the three node values are the cochains $e_a\wedge e_J,e_b\wedge e_J,e_c\wedge e_J$. Proposition~\ref{prop:two-supported} supplies a centre
\[
v=(e_a+e_c)\wedge e_I\wedge e_j.
\]
For an $(n+1)$-element set $U$, define $\delta_U(e_A)=1$ if the $n$-set $A$ is contained in $U$, and zero otherwise. On $(\sum e_i)\wedge e_B$ it counts zero or two faces, so it descends to~\eqref{eq:two-arrangement}.

The quotient row is exactly $\phi_I=\delta_{T\cup I}$. To check its normalization, split $W=A_T\oplus B_T$, where $A_T$ is spanned by $e_a+e_c,e_b+e_c$ and $B_T$ by the outside labels. This is a direct sum because three is invertible. The constant and node pieces are $\bigwedge^nB_T$ and $A_T\wedge\bigwedge^{n-1}B_T$, both annihilated by $\delta_{T\cup I}$. On
\[
(e_a+e_c)\wedge(e_b+e_c)\wedge e_{I'}
\]
its value is zero unless $I'=I$, and then is $3=1$. Refinement sends this vector to the sum of the three triangle-edge terms, three times the cycle generator at $I$, and to zero at all other elliptic indices. Hence the actual image contains
\begin{equation}\label{eq:two-marked-shear}
1+\bigl[(e_a+e_c)\wedge e_I\wedge e_j\bigr]\delta_{T\cup I}
\end{equation}
for all indicated labels. No splitting of~\eqref{eq:two-extension} was used to choose the row.

\section{The full characteristic-two groups}\label{sec:higher-two-group}
This section generates the entire covector stabilizer from the marked shears and identifies the parity in which that stabilizer enlarges to the full linear group.

Choose $e_1,\ldots,e_{d-1}$ as a basis of $W$, so $e_d=\sum_{r<d}e_r$. Let $v_A=e_A$, for $n$-subsets $A\subset\{1,\ldots,d-1\}$, and put $\ell=\sum_Av_A^*$. If $R$ is an $(n+1)$-subset, let $\mathcal Q(R)$ be its $n$-faces. For distinct $A,B,C\in\mathcal Q(R)$, write $R\setminus A=\{j\}$, $R\setminus B=\{a\}$, and $R\setminus C=\{c\}$. Apply~\eqref{eq:two-marked-shear} with $T=\{d,a,c\}$ and $I=A\setminus\{a,c\}$. Its row is $v_A^*$, since $A$ is the only $n$-face of $A\cup\{d\}$ not containing $d$, and its centre is $v_B+v_C$. Thus the image contains
\begin{equation}\label{eq:two-triple-operation}
1+(v_B+v_C)v_A^*
\end{equation}
for all triples of distinct coordinate vectors in each clique $\mathcal Q(R)$.

\begin{lemma}\label{lem:two-triple-stabilizer}
On a vector space with basis $f_1,\ldots,f_s$, $s\ge3$, the transformations
\[
1+(f_b+f_c)f_a^*,\qquad a,b,c\text{ distinct},
\]
generate the full stabilizer of $\ell_s=\sum_i f_i^*$ over $\mathbf F_2$.
\end{lemma}
\begin{proof}
All generators fix $\ell_s$. Put $K=\ker\ell_s$, with basis $g_i=f_i+f_s$ for $i<s$, and use $w=f_s$ as complement. The operations with distinct $a,b<s$ and $c=s$ act by $g_a\mapsto g_a+g_b$ and fix $w$. They generate $\mathrm{GL}(K)$, since over $\mathbf F_2$ the determinant group is trivial. An operation with $a=s$ and distinct $b,c<s$ takes $w$ to $w+g_b+g_c$. Cancel its restriction to $K$ using this copy of $\mathrm{GL}(K)$. The resulting nonzero translation, and its conjugates, give every map $w\mapsto w+z$ for $z\in K$. These translations together with $\mathrm{GL}(K)$ are exactly $\operatorname{Stab}(\ell_s)$.
\end{proof}

Each clique therefore contains all its coordinate permutations, fixing the other coordinates. The graph of $n$-subsets joined when they differ by one label is connected, and each adjacent pair lies in one clique. Its edge transpositions give every permutation of the full coordinate basis. Conjugating one transformation~\eqref{eq:two-triple-operation} by these permutations supplies every triple operation on that basis. The lemma gives
\begin{equation}\label{eq:two-stabilizer-lower}
\operatorname{Stab}(\ell)\subseteq G.
\end{equation}

If $d\equiv n+1\pmod2$, the canonical covector of Section~\ref{sec:pn-6} is nonzero at this ordinary arrangement and invariant under permutations of the hyperplanes. The subgroup $S_{d-1}$ acts transitively on the coordinate $n$-subsets, so all its coefficients are equal. Over $\mathbf F_2$, nonzero means they are all one. Thus it is $\ell$, and the canonical upper bound together with~\eqref{eq:two-stabilizer-lower} proves equality.

If $d\equiv n\pmod2$, let $\tau$ interchange labels $r$ and $d$. For a basis vector $v_A$ not containing $r$, $\ell(\tau v_A)=1$. If $r\in A$ and $A=B\cup\{r\}$, then
\[
\tau v_A=e_B\wedge e_d=
\sum_{j<d,\ j\notin B}e_B\wedge e_j.
\]
There are $d-n$ terms, an even number, so $\ell(\tau v_A)=0$. Both types of $A$ occur. Hence $\tau$ moves $\ell$.

A subgroup containing $\operatorname{Stab}(\ell)$ and an element moving $\ell$ is all of $\mathrm{GL}(V)$ in this situation. Indeed the stabilizer acts transitively on the other nonzero covectors: use $\mathrm{GL}(\ker\ell)$ on their nonzero restriction and its translation subgroup to adjust the complementary value. The group consequently contains a conjugate stabilizer of a coordinate row $f_1^*$. That stabilizer contains every $1+f_bf_a^*$ with $b\ne1$ and $a\ne b$. All coordinate permutations are already present, so conjugation gives every off-diagonal elementary matrix. They generate $\mathrm{GL}(V)$ over $\mathbf F_2$. This proves the geometric claims of Theorem~\ref{thm:higher-images} in characteristic two. The arithmetic image contains the geometric one, and in the stabilizer case the canonical covector descends over the perfect ground field. The arithmetic images are therefore equal to the asserted geometric images.


\section{The full classification and method}
The results above determine the residue threshold and every finite monodromy image for projective hypersurfaces. Theorem~\ref{thm:threshold} gives the exact threshold in every characteristic and dimension. Theorems~\ref{thm:images}, \ref{thm:plane-images}, and~\ref{thm:higher-images} give the zero- and rank-one cases, every plane group, and every higher-dimensional group. The characteristic-two parity obstruction is exact: the image is the full canonical-covector stabilizer, and the opposite parity has the full linear group.

For example, a plane quartic in characteristic two has $h=3$ and image $\mathrm{GL}_3(\mathbf F_2)$, whereas a plane quintic has $h=6$ and image $\mathbf F_2^5\rtimes\mathrm{GL}_5(\mathbf F_2)$. In $\mathbf P^3$, degree five has image $\mathrm{GL}_4(\mathbf F_2)$, and degree six has image $\mathbf F_2^9\rtimes\mathrm{GL}_9(\mathbf F_2)$. Every one of these families realizes both residues in every nonempty parameter open after finite extension. The image calculations distinguish the complete finite representation even when the trace conclusion is the same.

The original computational supplement \emph{Projective Frobenius formula checks}, at the \href{\workbenchurl}{workbench}, checks point counts, semilinear matrix composition, the differential identity over $\mathbf F_2$ and $\mathbf F_4$, the rank-one polynomial gcd, and distinct coefficient positions. The additional original supplement \emph{Higher hypersurface matrix checks}, at the same workbench, verifies finite exterior-coordinate overlap and spanning calculations and the explicit characteristic-three minors. The original supplement \emph{Small-characteristic hypersurface checks} verifies the elliptic leading coefficients, marked cochain identities, modular norms, elementary matrix words, and selected small finite groups. These finite computations corroborate formulas; the universal monodromy conclusions rely on the proofs above.

\paragraph{Sources and method.}
The mathematics and text were produced by GPT-6 Astra (ChatGPT Pro mode). The residue theorem, joint cubic-boundary image, determinant consequence, plane groups, and higher-dimensional image theorems have completed separate statement-led same-model reconstructions and full comparisons at their recorded input scopes. The characteristic-three first-degree proof has two distinct coefficient charts; the higher-prime proof has different extension and elementary-matrix arguments. The higher small-characteristic proofs have separate complete reconstructions, with the earlier cubic-boundary statements used as explicit inputs. The appendix gives complete targeted proofs of the additional restriction and elementary group lemmas; these receive no extra independent-main-proof credit. Related prior exposure is recorded; no globally pristine or different-model independence is claimed. Full mathematical and newcomer readings, exact primary-source comparisons, and fresh source-only builds with page inspection are recorded at their actual scopes. Reasoning-effort metadata is not exposed. No human has checked the mathematics. Finite computations corroborate formulas and small groups; the universal conclusions rely on the written proofs. Our text and original code are dedicated to the public domain under CC0~1.0. The editable paper and programs are at \url{\workbenchurl}; the permanent address of its records is \url{\recordurl}.

\appendix
\section{Restriction geometry and elementary group lemmas}\label{app:auxiliaries}
This appendix records useful auxiliary statements about the parameter spaces and finite representations. The full image theorems above already imply irreducibility, their special linear lower bounds, and their determinant consequences. The statements below give additional restriction and generating-set information without changing those theorems.

For comparison, the symmetric-group module at an arrangement has a classical description. Let $P=\mathbf F_p^d$, $u=\sum e_i$, $U=\ker(\sum:P\to\mathbf F_p)$, $W=P/\langle u\rangle$, and let $D$ be the image of $U$ in $W$. For odd $p$ and $1\le r\le d-2$, the module $\bigwedge^rW$ is simple when $p\nmid d$. When $p\mid d$, its unique proper nonzero submodule is $\bigwedge^rD$, and its simple quotient is $\bigwedge^{r-1}D$. This is the dual of the augmentation-hook description in \citet[Theorem~2, p.~142; Propositions~1 and~4, pp.~136--138]{Peel}. In particular, no semisimplicity of the entire symmetric-group representation is asserted in the divisible case.

\subsection{Three modular group lemmas}\label{app:modular}
These elementary arguments describe cocycles, normal transvection blocks and a reflection obstruction. They require no classification of finite linear groups.

\begin{lemma}\label{lem:sl3-cocycles}
For an odd prime $p$ and the natural module $Q=\mathbf F_p^3$,
\[
H^1(\mathrm{SL}_3(\mathbf F_p),Q)
=H^1(\mathrm{SL}_3(\mathbf F_p),Q^*)=0.
\]
\end{lemma}
\begin{proof}
Let $c$ be a cocycle, so $c(gh)=c(g)+gc(h)$. Let $D_0$ be the diagonal sign subgroup of order four. Averaging over $D_0$, whose order is invertible, permits subtraction of a coboundary so that $c$ vanishes on that subgroup. More explicitly, if $v=|D_0|^{-1}\sum_{s\in D_0}c(s)$, then $c(t)=v-tv$ for $t\in D_0$; subtract this coboundary.

Write $t_{ij}(a)=1+aE_{ij}$. This element commutes with the sign matrix that is negative on coordinates $i,j$. The cocycle relation therefore makes $c(t_{ij}(a))$ fixed by that sign matrix, hence supported on the remaining coordinate. For
\[
x=t_{12}(a),\quad y=t_{23}(b),\quad
c(x)=Ae_3,\quad c(y)=Be_1,
\]
the cocycle identity and $c(g^{-1})=-g^{-1}c(g)$ give
\[
c([x,y])=-bAe_2-abAe_1.
\]
But $[x,y]=t_{13}(ab)$ has cocycle value supported on $e_2$. Taking $a,b\ne0$ forces $A=0$. Relabelling the coordinates proves vanishing on every elementary matrix. These matrices generate $\mathrm{SL}_3(\mathbf F_p)$, so the normalized cocycle is zero. The inverse-transpose automorphism identifies the dual calculation with the natural one.
\end{proof}

\begin{lemma}\label{lem:normal-transvection-blocks}
Let a finite subgroup $G\subseteq\mathrm{GL}(V)$ act irreducibly over $\mathbf F_p$ and contain a transvection. Let $N$ be its normal subgroup generated by all transvections in $G$. Then $V$ is a direct sum of pairwise nonisomorphic simple $N$-modules, permuted transitively by $G$. Each transvection has its image line in one summand and fixes every other summand pointwise.
\end{lemma}
\begin{proof}
The $N$-socle, the sum of all simple $N$-submodules, is nonzero. Normality makes it $G$-stable, so it is all of $V$. Thus restriction to $N$ is semisimple. The canonical isotypic summands are permuted transitively by $G$, since the sum over any orbit is $G$-stable.

On an isotypic summand $S^a$, every $x\in N$ acts as the direct sum of $a$ copies of its action on $S$. Its rank $\operatorname{rank}(x-1)$ there is therefore a multiple of $a$. A transvection is nontrivial and has total rank one, so on its unique nontrivial isotypic summand one has $a=1$. Transitivity gives multiplicity one on every summand. The rank-one condition then also shows that a transvection acts nontrivially on just one summand. This proves every assertion without applying Maschke's theorem.
\end{proof}

\begin{lemma}\label{lem:symplectic-reflection-obstruction}
Suppose $p$ is odd and an irreducible normal subgroup $N\triangleleft G\subseteq\mathrm{GL}(V)$ is the natural symplectic group, with $\dim V\ge3$. Then $G$ contains no rank-one reflection, meaning an element with eigenvalue $-1$ once and eigenvalue $1$ on a hyperplane.
\end{lemma}
\begin{proof}
Fix the invariant nondegenerate alternating form $B$. Its invariant alternating forms are scalar multiples of $B$. Indeed, writing another such form as $B(Lv,w)$ makes $L$ commute with $N$. Commuting with the symplectic transvection of each nonzero centre $v$ preserves its image line $\langle v\rangle$. Hence $L$ preserves every line and is scalar. Normality of $N$ now implies that each $g\in G$ is a similitude of $B$.

Suppose a reflection $R$ has positive hyperplane $H$ and multiplier $c$. If $c\ne1$, then $B|_{H\times H}=0$. But an isotropic subspace has dimension at most $\dim V/2$, whereas $\dim H=\dim V-1$; the symplectic dimension is even and at least four. If $c=1$, the negative eigenline is perpendicular to $H$, since $p$ is odd. It is also perpendicular to itself because $B$ is alternating, so it lies in the radical of $B$. Both cases are impossible.
\end{proof}

\subsection{Projection and prescribed restrictions}\label{app:projection}
The next statements preserve the full finite image while prescribing a geometric generic restriction of an equation.

\begin{lemma}\label{lem:projection-connected}
Let $k$ be algebraically closed, let $U\subset\mathbf P^N_k$ be nonempty open, and let $T\to U$ be a connected finite \textup{\'etale} cover. Let $f:\mathbf P^N\dashrightarrow\mathbf P^M$ be a linear projection whose centre $Z$ meets $U$. The inverse image in $T$ of the geometric generic fibre of $f|_{U\setminus Z}$ is connected. Consequently a finite local system on $U$ retains its full image on that geometric generic fibre.
\end{lemma}
\begin{proof}
The scheme $T$ is smooth and integral, by connectedness and normality; see \citet[Lemma~58.11.2]{StacksNormal}. Pull the blowup of $\mathbf P^N$ along $Z$ back to $T$, obtaining a smooth integral scheme $T'$ with a morphism to $\mathbf P^M$. Choose $z\in Z\cap U$ and a point of $T$ above it. Both are $k$-rational. The exceptional fibre at this point is $\mathbf P^M$, and the resolved projection maps this fibre identically to the target. Thus $T'\to\mathbf P^M$ has a section.

Put $K=k(\mathbf P^M)$ and let $L$ be the algebraic closure of $K$ in $k(T')$. This is a finite extension. Let $C$ be the normalization of $\mathbf P^M$ in $L$. It is finite, since a finite-type scheme over a field is Nagata. The morphism $T'\to\mathbf P^M$ factors through $C$: an element of $L$ integral over a coordinate ring of the target is integral over every normal local ring of $T'$ above that affine open, and hence belongs to it. These local maps glue. This is the normalization factorization of \citet[Lemma~29.54.4]{StacksNormalization}.

The exceptional section induces a section $\mathbf P^M\to C$. Its image is closed, because $C$ is finite over the target, and has the same dimension as the integral scheme $C$. It is therefore all of $C$; reducedness then makes the section an isomorphism. Consequently $L=K$.

Away from the exceptional divisor, the map is \textup{\'etale} over an open of the affine-space fibres of the linear projection. Its generic fibre is smooth. The equality $L=K$ implies geometric irreducibility of that generic fibre by \citet[Lemma~33.8.6]{StacksGeometric}. Smoothness gives geometric reducedness, hence geometric integrality. Removing the exceptional divisor leaves a nonempty geometrically integral open. This is the required inverse image in $T$.

Finally, a connected component of the finite frame cover for a local system is connected after this restriction. Its deck group is its finite image, so that image is preserved.
\end{proof}

For $j\ge2$, let $W_{j,e}$ denote the Hasse--Witt invertibility open in the full projective space of degree-$e$ equations in $\mathbf P^j$, including singular equations. Let $V_{j,e}$ be the local system of Frobenius-fixed vectors in $H^{j-1}(\mathcal O_Y)$, and let $G_{j,e}$ be its geometric image. Sections~\ref{sec:pn-1}--\ref{sec:pn-2} give its constant rank $\binom{e-1}{j}$ and its finite \textup{\'etale} construction.

\begin{proposition}\label{prop:restriction-product}
Let $k$ be algebraically closed of characteristic $p$, let $n\ge3$ and $e\ge n+1$, and fix a hyperplane $H\subset\mathbf P^n$. On the open parameter space where both a degree-$e$ equation and its restriction to $H$ have invertible Hasse--Witt maps, the joint geometric image on $V_{n,e}$ and $V_{n-1,e}$ is the product of their full geometric images. In particular, $V_{n,e}$ retains its full image when the restriction to $H$ is prescribed to be geometric generic.
\end{proposition}
\begin{proof}
Restriction of equations is a linear projection between their projective parameter spaces. Its centre consists of equations divisible by the equation of $H$. A general-position union of $e$ hyperplanes, one of them $H$, lies in that centre and in $W_{n,e}$: its intersection complex has constant coefficients and Frobenius is invertible. Hence Lemma~\ref{lem:projection-connected} applies to the finite frame cover of $V_{n,e}$.

Away from its centre, restriction is an affine-space bundle. The additional invertibility opens have nonempty intersection: the space of nonzero restrictions is dominant, and each invertibility condition is a nonempty open. Its geometric generic fibre is therefore a nonempty open of affine space, hence geometrically integral. Pullback along this map preserves connected finite \textup{\'etale} covers of the target, so the image on $V_{n-1,e}$ is full. On the geometric generic fibre, the image on $V_{n,e}$ is full by the preceding lemma and the restricted representation is constant. The joint image consequently contains $G_{n,e}\times1$ and surjects onto $G_{n-1,e}$. These two properties imply that it is the product.
\end{proof}

\begin{lemma}\label{lem:elliptic-constant-section}
Suppose $p$ is odd and $k$ is algebraically closed. There is a family of plane cubics over a strictly henselian smooth trait over $k$ whose special fibre is smooth supersingular, whose Hasse invariant has valuation one, and whose restriction to a fixed line is the constant binary form
\[
g=x_0x_1(x_0-x_1).
\]
There is also a smooth ordinary plane cubic with this exact restriction.
\end{lemma}
\begin{proof}
Section~\ref{sec:pn-5} proves that the Legendre Hasse polynomial has a simple root away from $0$ and $1$. Take the strict henselization of the parameter line at such a root. The resulting plane-cubic family is smooth, its special fibre is supersingular, and its Hasse invariant has valuation one.

Choose a line transverse to the special cubic. Its intersection with the family is finite \textup{\'etale} of degree three, so over the strictly henselian trait it splits into three disjoint sections. A regular automorphism of that projective line takes these ordered sections to $0,1,\infty$. To see regularity explicitly, represent the first two sections by primitive vectors over the local ring; disjointness makes them a basis. The third has two unit coordinates in this basis, and diagonal rescaling sends it to $(1,1)$. The resulting invertible two-by-two matrix extends to an invertible plane-coordinate matrix by adjoining one diagonal entry.

Apply this coordinate change to the entire family. Its restricted binary cubic now has the zeros of $g$ and differs from $g$ by a unit, since those three zeros remain distinct. Rescale the defining equation by the inverse unit. The restriction is exactly $g$. Coordinate changes and unit rescaling preserve smoothness, supersingularity and the valuation of the Hasse section: they change its trivialization only by a unit. The ordinary assertion follows from the prescribed-squarefree-restriction construction in Section~\ref{sec:plane-cubic-restriction}.
\end{proof}

\subsection{An explicit characteristic-two generating set}\label{app:coordinate-generators}
This calculation identifies the subgroup generated by a specific family of coordinate swaps, giving more information than their membership in the full monodromy image.

Let $n\ge2$, $d\ge n+2$, $m=d-1$, and
\[
W=\mathbf F_2^d/\mathbf F_2\mathbf1,
\qquad V=\bigwedge^nW,
\qquad h=\binom mn.
\]
Choose the classes $e_1,\ldots,e_m$ as a basis, so $e_0=\sum_{i=1}^m e_i$. Write $b_A=\bigwedge_{i\in A}e_i$ for its exterior basis and $\ell=\sum_A b_A^*$. For an $(n-1)$-subset $I\subset\{2,\ldots,m\}$ put
\[
w_I=\sum_{a\in\{2,\ldots,m\}\setminus I}b_{I\cup\{a\}},
\qquad R_I=1+w_Ib_{I\cup\{1\}}^*.
\]
The $R_I$ are nontrivial commuting involutions: their rows are distinct basis rows containing label $1$, whereas all their centres avoid that label. A product is the identity only when every chosen factor is absent, by inspection of its distinct input columns. They thus form an independent elementary abelian subgroup. In the stabilizer parity their centres have even weight, so these operators belong to the image already determined by the main theorems.

\begin{proposition}\label{prop:coordinate-generators}
Let $G_0$ be generated by the permutation action of $S_d$ on $V$ and all the $R_I$. Then
\[
G_0=\begin{cases}
\mathrm{GL}_h(\mathbf F_2),&d-n-1\text{ odd},\\
\operatorname{Stab}_{\mathrm{GL}_h(\mathbf F_2)}(\ell),&d-n-1\text{ even}.
\end{cases}
\]
In the second case the group is $\mathbf F_2^{h-1}\rtimes\mathrm{GL}_{h-1}(\mathbf F_2)$.
\end{proposition}
\begin{proof}
Set $l=d-n-1=m-n$. Conjugating by the subgroup $S_m$ fixing label $0$ gives, for each $(n-1)$-set $I$ and each $j\notin I$, the involution $T_j^I$ which sends $b_{I\cup\{j\}}$ to
\[
S_I=\sum_{a\in\{1,\ldots,m\}\setminus I}b_{I\cup\{a\}}
\]
and fixes all other basis vectors. It also sends $S_I$ back to $b_{I\cup\{j\}}$. For distinct $j,t\notin I$, the word $T_j^IT_t^IT_j^I$ therefore exchanges just $b_{I\cup\{j\}}$ and $b_{I\cup\{t\}}$. The graph of $n$-subsets joined when they differ by one label is connected: replace successively a label outside the target subset by one in the target. Its edge transpositions generate every permutation of the coordinate basis. Thus $G_0$ contains $S_h$.

Each $R_I$ has a centre of weight $l$ disjoint from its input coordinate. Coordinate conjugation now supplies every transformation $1+w b_A^*$ with any such centre. If $l=1$, these are all elementary matrices, proving the first case. Suppose $l\ge2$. Since $h=\binom mn\ge m=l+n\ge l+2$, one can choose two $l$-element supports, both avoiding a fixed input coordinate, with intersection of size $l-1$ and prescribed distinct remaining coordinates $B,C$. Multiplying the same-row transformations gives
\[
1+(b_B+b_C)b_A^*.
\]
Hence all triple operations of Lemma~\ref{lem:two-triple-stabilizer} occur.

If $l$ is odd, add weight-two centres within an $l$-element centre to cancel pairs until one coordinate remains. Same-row transformations multiply by adding their centres. Coordinate permutations then give every elementary matrix, so $G_0=\mathrm{GL}_h(\mathbf F_2)$.

If $l$ is even, the $R_I$ and the coordinate subgroup $S_m$ preserve $\ell$. The arrangement transposition $(01)$ is exactly $\prod_I R_I$: on a basis vector not containing label $1$ both sides are the identity, and on $b_{I\cup\{1\}}$ both replace $e_1$ by $e_0=\sum e_i$. Since $S_m$ and $(01)$ generate $S_d$, every generator of $G_0$ preserves $\ell$. Conversely the triple operations generate its full stabilizer by Lemma~\ref{lem:two-triple-stabilizer}, whose elementary proof includes both the linear action on $\ker\ell$ and all translations of a complementary vector. This proves the second equality and its semidirect-product description.
\end{proof}


\begin{thebibliography}{99}
\bibitem[Grothendieck(n.d.)]{Grothendieck}
Grothendieck, A. (n.d.). \emph{Notes on crystalline cohomology}, cote~4, \S8.6, archive sheet~11, PDF image~12. Archives Alexander Grothendieck, Universit\'e de Montpellier. \url{https://grothendieck.umontpellier.fr/4.pdf}.
\bibitem[Adolphson and Sperber(2016)]{AdolphsonSperber}
Adolphson, A., \& Sperber, S. (2016). $A$-hypergeometric series and the Hasse--Witt matrix of a hypersurface. \emph{Finite Fields and Their Applications, 41}, 55--63. \url{https://doi.org/10.1016/j.ffa.2016.05.001}. Author version: \url{https://arxiv.org/abs/1601.05127}; Theorem~2.3 and Proposition~2.11, pp.~2--3.
\bibitem[Koblitz(1975)]{Koblitz}
Koblitz, N. (1975). $p$-adic variation of the zeta-function over families of varieties defined over finite fields. \emph{Compositio Mathematica, 31}(2), 119--218. \url{https://www.numdam.org/item/CM_1975__31_2_119_0/}. Theorem~3, p.~142, with proof pp.~142--144; the hypersurface result is credited there to Miller, reference~[40].
\bibitem[Huang et~al.(2023)]{HLYY}
Huang, A., Lian, B., Yau, S.-T., \& Yu, C. (2023).
Hasse--Witt matrices, unit roots and period integrals.
\emph{Mathematische Annalen, 387}, 145--173.
\url{https://doi.org/10.1007/s00208-022-02464-y}.
Author version: \url{https://arxiv.org/abs/1801.01189};
Corollaries 2.3 and 2.8, pp.\ 8--11; Theorem 5.3, pp.\ 19--22.

\bibitem[OpenAI(2026)]{OpenAI}
OpenAI. (2026, October~7). \emph{Milne's rationality conjecture for abelian varieties} [Preprint]. \url{https://github.com/openai/math/tree/main/preprints/Milnes-rationality-conjecture-for-abelian-varieties-October-7-2026}. Theorem~1.1, pp.~1--2; context only.
\bibitem[Katz(1972)]{Katz1972}
Katz, N. M. (1972). Algebraic solutions of differential equations ($p$-curvature and the Hodge filtration). \emph{Inventiones Mathematicae, 18}(1--2), 1--118. \url{https://doi.org/10.1007/BF01389714}. Author PDF: \url{https://web.math.princeton.edu/~nmk/old/algsoln.pdf}. \S2.3.7, pp.~34--36, especially Algorithm~2.3.7.14 and formula~2.3.7.16.
\bibitem[Katz(1973)]{Katz}
Katz, N. M. (1973). Une formule de congruence pour la fonction $\zeta$. In P. Deligne and N. Katz (Eds.), \emph{Groupes de monodromie en g\'eom\'etrie alg\'ebrique, II}, Expos\'e~XXII (Lecture Notes in Mathematics, Vol.~340, pp.~401--438). Springer. \url{https://web.math.princeton.edu/~nmk/old/congzeta.pdf}.
\bibitem[Lang and Weil(1954)]{LangWeil}
Lang, S., \& Weil, A. (1954). Number of points of varieties in finite fields. \emph{American Journal of Mathematics, 76}(4), 819--827. \url{https://doi.org/10.2307/2372655}.
\bibitem[Peel(1971)]{Peel}
Peel, M. H. (1971). Hook representations of the symmetric groups. \emph{Glasgow Mathematical Journal, 12}(2), 136--149. \url{https://doi.org/10.1017/S0017089500001245}. Theorem~2, p.~142; Propositions~1 and~4, pp.~136--138.
\bibitem[Poonen(2017)]{Poonen}
Poonen, B. (2017). \emph{Rational points on varieties} (Graduate Studies in Mathematics, Vol.~186). American Mathematical Society. Theorem~7.7.1(iii), pp.~222--223. \url{https://math.mit.edu/~poonen/papers/Qpoints.pdf}.
\bibitem[The Stacks Project Authors(2026a)]{StacksProjective}
The Stacks Project Authors. (2026a). \emph{Cohomology of projective space}, Section~30.8, Lemmas~30.8.1--30.8.2, Tag~01XS. \url{https://stacks.math.columbia.edu/tag/01XS}.
\bibitem[The Stacks Project Authors(2026b)]{StacksAS}
The Stacks Project Authors. (2026b). \emph{The Artin--Schreier sequence}, Section~59.63, Tag~0A3J. \url{https://stacks.math.columbia.edu/tag/0A3J}.
\bibitem[The Stacks Project Authors(2026c)]{StacksComparison}
The Stacks Project Authors. (2026c). \emph{Cohomology of quasi-coherent sheaves}, Theorem~59.22.4, Tag~03P2. \url{https://stacks.math.columbia.edu/tag/03P2}.
\bibitem[The Stacks Project Authors(2026d)]{StacksProperAS}
The Stacks Project Authors. (2026d). \emph{Lemma 59.63.6}, Tag~0A3P. \url{https://stacks.math.columbia.edu/tag/0A3P}.
\bibitem[The Stacks Project Authors(2026e)]{StacksNormal}
The Stacks Project Authors. (2026e). \emph{Fundamental groups of normal schemes}, Lemma~58.11.2, Tag~0BQL. \url{https://stacks.math.columbia.edu/tag/0BQL}.
\bibitem[The Stacks Project Authors(2026f)]{StacksPurity}
The Stacks Project Authors. (2026f). \emph{Purity of branch locus}, Lemma~58.21.4, Tag~0BMB. \url{https://stacks.math.columbia.edu/tag/0BMB}. The normalizations used here are finite and dominate the regular base; no vertical component occurs.
\bibitem[The Stacks Project Authors(2026g)]{StacksRH}
The Stacks Project Authors. (2026g). \emph{Riemann--Hurwitz}, Section~53.12, especially Lemma~53.12.2, Tag~0C1B. \url{https://stacks.math.columbia.edu/tag/0C1B}.
\bibitem[The Stacks Project Authors(2026h)]{StacksBaseChange}
The Stacks Project Authors. (2026h). \emph{Cohomology and base change, III}, Lemma~30.22.1, Tag~07VK. \url{https://stacks.math.columbia.edu/tag/07VK}.
\bibitem[The Stacks Project Authors(2026i)]{StacksPlane}
The Stacks Project Authors. (2026i). \emph{Plane curves}, Section~53.9, Lemma~53.9.3, Tag~0BYA. \url{https://stacks.math.columbia.edu/tag/0BYA}.
\bibitem[The Stacks Project Authors(2026j)]{StacksNormalization}
The Stacks Project Authors. (2026j). \emph{Lemma 29.54.4}, Tag~035I. \url{https://stacks.math.columbia.edu/tag/035I}.
\bibitem[The Stacks Project Authors(2026k)]{StacksGeometric}
The Stacks Project Authors. (2026k). \emph{Geometrically irreducible schemes}, Section~33.8, especially Lemma~33.8.6, Tag~0364. \url{https://stacks.math.columbia.edu/tag/0364}.
\end{thebibliography}
\end{document}
